\documentclass{beamer}
\usetheme{Madrid}
\usecolortheme{default}
\usepackage[utf8]{inputenc}
\usepackage{amsmath}
\usepackage{amsfonts}
\usepackage{amssymb}
\usepackage{tikz}


\title{CRISC Exam - Hard Questions with Explanations}
\subtitle{From   CRISC GPTs}
\author{Compiled from ISACA CRISC GPT APIs}
\date{\today}

\begin{document}
	
	\begin{frame}
		\titlepage
	\end{frame}
	
	\begin{frame}{Outline}
		\tableofcontents
	\end{frame}
	
	% ================ SECTION 1: DOMAIN 1 - IT RISK IDENTIFICATION ================
	
	\section{Domain 1: IT Risk Identification}
	
	\begin{frame}{Q1: Resilient Business Processes}
		\begin{block}{Question}
			Which business requirement BEST relates to the need for resilient business and information systems processes?
		\end{block}
		\begin{enumerate}[(A)]
			\item Effectiveness
			\item Confidentiality
			\item Integrity
			\item Availability
		\end{enumerate}
	\end{frame}
	
	\begin{frame}{Q1: Explanation}
		\begin{block}{Correct Answer: D}
			Availability relates to information being available when required by the business process now and in the future. Resilience is the ability to provide and maintain an acceptable level of service during disasters or when facing operational challenges.
		\end{block}
		\begin{block}{Option Analysis}
			\begin{itemize}
				\item \textbf{A - Incorrect:} Effectiveness deals with information being relevant and delivered timely
				\item \textbf{B - Incorrect:} Confidentiality deals with protection from unauthorized disclosure
				\item \textbf{C - Incorrect:} Integrity relates to accuracy and completeness
				\item \textbf{D - Correct:} Availability is directly linked to resilience
			\end{itemize}
		\end{block}
	\end{frame}
	
	
	\begin{frame}{Q1: Resilient Business Processes} \begin{block}{Question} Which business requirement BEST relates to the need for resilient business and information systems processes? \end{block} \begin{enumerate}[(A)] \item Effectiveness \item Confidentiality \item Integrity \item Availability \end{enumerate} \begin{block}{Memory Aid} Confidentiality = Prevent Disclosure Integrity = Prevent Unauthorized Change Availability = Keep Systems Running Resilience = Availability During Disruption \end{block} \end{frame} \begin{frame}{Q1: Explanation} \footnotesize \begin{block}{Correct Answer: D - Availability} Resilience is the ability of systems and business processes to continue operating or recover quickly following a disruption. This objective is most closely aligned with availability because information and services must remain accessible when required. \end{block} \begin{block}{Option Analysis} \textbf{A - Incorrect: Effectiveness} Effectiveness focuses on whether information supports business objectives, decision-making, and operational needs. While an effective process may contribute to resilience, effectiveness does not specifically address continuity during disruptions. \vspace{0.15cm} \textbf{B - Incorrect: Confidentiality} Confidentiality ensures information is protected from unauthorized disclosure. A highly confidential system can still be unavailable during an outage, making confidentiality different from resilience. \vspace{0.15cm} \textbf{C - Incorrect: Integrity} Integrity ensures information remains accurate, complete, and authorized. Although integrity is critical, resilience focuses on maintaining service availability rather than preventing unauthorized modification. \vspace{0.15cm} \textbf{D - Correct: Availability} Availability ensures systems, data, and services remain accessible when needed. Disaster recovery, business continuity, fault tolerance, redundancy, and resiliency initiatives are all primarily intended to support availability. \end{block}  \end{frame}
	
	\begin{frame}{Q2: Value of Risk Register}
		\begin{block}{Question}
			Which statement BEST describes the value of a risk register?
		\end{block}
		\begin{enumerate}[(A)]
			\item It captures the risk inventory
			\item It drives the risk response plan
			\item It is a risk reporting tool
			\item It lists internal risk and external risk
		\end{enumerate}
		\begin{block}{Memory Aid} Risk Assessment = Identify Risks Risk Register = Record Risks Risk Owner = Own the Risk Risk Register = Prioritize Risks Risk Response Plan = Treat Risks Reports/Dashboards = Communicate Risks \vspace{0.15cm} Think: "No Register = No Prioritization" "No Prioritization = No Risk Response" \end{block}
	\end{frame}
	
\begin{frame}{Q2: Explanation} \footnotesize \begin{block}{Correct Answer: B - It Drives the Risk Response Plan} A risk register is the central repository for risk information. It documents risk descriptions, likelihood, impact, ownership, treatment strategies, and residual risk. Its greatest value is enabling management to prioritize and respond to risks in a consistent manner. \end{block} \begin{block}{Option Analysis} \textbf{A - Incorrect: It Captures the Risk Inventory} A risk register does maintain an inventory of identified risks. However, simply listing risks provides little value unless the information is used to determine ownership, prioritization, treatment, and monitoring activities. \vspace{0.15cm} \textbf{B - Correct: It Drives the Risk Response Plan} Risk response decisions are based on the information contained within the risk register. Management uses risk ratings, ownership assignments, treatment status, and residual risk information to determine how risks should be addressed. \vspace{0.15cm} \textbf{C - Incorrect: It Is a Risk Reporting Tool} Risk reports and dashboards may be generated from the risk register, but the register itself is primarily a repository and management tool rather than a reporting mechanism. \vspace{0.15cm} \textbf{D - Incorrect: It Lists Internal Risk and External Risk} A risk register may contain both internal and external risks, but this is merely one attribute of the information stored. The value of the register comes from supporting risk-management decisions, not categorizing risks. \end{block}  \end{frame}
	
	\begin{frame}{Q3: New Law Impact on HR System}
		\begin{block}{Question}
			Shortly after annual policy review, a new law may affect security requirements for HR system. Risk practitioner should:
		\end{block}
		\begin{enumerate}[(A)]
			\item Analyze how the law may affect the enterprise
			\item Ensure adjustments during next review cycle
			\item Initiate an ad hoc revision of corporate policy
			\item Notify system custodian to implement changes
		\end{enumerate}
	\end{frame}
	
	\begin{frame}{Q3: Explanation}
		\begin{block}{Correct Answer: A}
			Assessing how the law may affect the enterprise is the best course of action. The analysis must also determine whether existing controls already address the new requirements.
		\end{block}
		\begin{block}{Option Analysis}
			\begin{itemize}
				\item \textbf{A - Correct:} Analysis is the first step
				\item \textbf{B - Incorrect:} Significant changes should trigger ad hoc reassessment
				\item \textbf{C - Incorrect:} Policy should be developed systematically
				\item \textbf{D - Incorrect:} Changes need approval by process owner
			\end{itemize}
		\end{block}
	\end{frame}
	
	
	\begin{frame}{Q3: Regulatory Change Impact} \begin{block}{Question} Shortly after the annual risk assessment is completed, a new privacy regulation is enacted that may affect several business processes involving employee information. Senior management asks the risk practitioner how the organization should respond. What should the risk practitioner do FIRST? \end{block} \begin{enumerate}[(A)] \item Assess the impact of the regulation on business processes and existing controls \item Initiate an immediate update to all security policies \item Direct system owners to implement new controls \item Wait until the next scheduled review cycle \end{enumerate} \begin{block}{CRISC Exam Tip} New Law $\rightarrow$ Impact Assessment $\rightarrow$ Gap Analysis $\rightarrow$ Policy/Control Updates $\rightarrow$ Implementation When asked what to do FIRST, choose the answer involving assessment, analysis, or understanding the situation. \end{block} \end{frame} \begin{frame}{Q3: Explanation} \footnotesize \begin{block}{Correct Answer: A} Before modifying policies, controls, or systems, management must first understand how the regulatory change affects the organization. The impact assessment determines whether changes are required and whether existing controls already satisfy the new requirements. \end{block} \begin{block}{Option Analysis} \textbf{A - Correct} Impact analysis is the appropriate first step. CRISC questions frequently test the principle that risk practitioners should assess before acting. Regulatory changes do not automatically require policy or control changes. \vspace{0.15cm} \textbf{B - Incorrect} Policies may eventually require revision, but updating policies before understanding the impact could create unnecessary work or introduce requirements that are not needed. \vspace{0.15cm} \textbf{C - Incorrect} Control implementation should occur only after management understands the regulatory requirements and determines which gaps actually exist. \vspace{0.15cm} \textbf{D - Incorrect} Significant regulatory changes may require action before the next review cycle. Deferring analysis could expose the organization to compliance risk. \end{block}  \end{frame}
	
	\begin{frame}{Q4: Weather Forecast System Priority}
		\begin{block}{Question}
			An information system that processes weather forecasts for public consumption is MOST likely to place its highest priority on:
		\end{block}
		\begin{enumerate}[(A)]
			\item Nonrepudiation
			\item Confidentiality
			\item Integrity
			\item Availability
		\end{enumerate}
	\end{frame}
	
	\begin{frame}{Q4: Explanation}
		\begin{block}{Correct Answer: C}
			A system that delivers weather forecasts is likely to place its highest priority on the integrity of the data. The risk practitioner should keep in mind that whether a forecast turns out to be accurate in its prediction is distinct from whether the data was accurately represented.
		\end{block}
		\begin{block}{Option Analysis}
			\begin{itemize}
				\item \textbf{A - Incorrect:} Nonrepudiation is unlikely to be important for weather forecasts
				\item \textbf{B - Incorrect:} Keeping data confidential would be at odds with public use
				\item \textbf{C - Correct:} Integrity of data is paramount
				\item \textbf{D - Incorrect:} Availability is important but integrity is higher priority
			\end{itemize}
		\end{block}
	\end{frame}
	
	
	\begin{frame}{Q4: Weather Forecast System Priority} \begin{block}{Question} An information system distributes weather forecasts to the public. Which information security objective should receive the HIGHEST priority? \end{block} \begin{enumerate}[(A)] \item Nonrepudiation \item Confidentiality \item Integrity \item Availability \end{enumerate} \begin{block}{CRISC Memory Aid} Weather Forecast = Integrity Emergency Services = Availability Medical Records = Confidentiality + Integrity Contracts/Digital Signatures = Nonrepudiation \end{block} \end{frame} \begin{frame}{Q4: Explanation} \footnotesize \begin{block}{Correct Answer: C - Integrity} The value of a weather forecasting system depends on the accuracy and reliability of the information being presented. If forecast data is altered, corrupted, or incorrectly displayed, the public could make poor decisions based on inaccurate information. Therefore, integrity is the most important objective. \end{block} \begin{block}{Option Analysis} \textbf{A - Incorrect: Nonrepudiation} Nonrepudiation provides proof that a specific party performed an action or submitted information. This requirement is far more important in areas such as contracts, financial transactions, and digital signatures than in public weather forecasting. \vspace{0.15cm} \textbf{B - Incorrect: Confidentiality} Weather forecasts are intended for public consumption. Restricting access to the information would not support the organization's mission. \vspace{0.15cm} \textbf{C - Correct: Integrity} Forecasts must be presented accurately and without unauthorized modification. Users depend on the correctness of the information, even if the prediction itself later proves inaccurate. Integrity focuses on ensuring the data has not been improperly altered. \vspace{0.15cm} \textbf{D - Incorrect: Availability} Availability is certainly important because users need access to forecasts when severe weather is approaching. However, an available system that provides incorrect information may create greater harm than a temporary outage. \end{block}  \end{frame}
	
	\begin{frame}{Q5: Best Risk Management View}
		\begin{block}{Question}
			Which provides the BEST view of risk management?
		\end{block}
		\begin{enumerate}[(A)]
			\item An interdisciplinary team
			\item A third-party risk assessment
			\item IT department risk assessment
			\item Internal compliance department
		\end{enumerate}
	\end{frame}
	
	\begin{frame}{Q5: Explanation}
		\begin{block}{Correct Answer: A}
			Having an interdisciplinary team contribute to risk management ensures that all areas are adequately considered and included in the risk assessment processes to support an enterprise view of risk.
		\end{block}
		\begin{block}{Option Analysis}
			\begin{itemize}
				\item \textbf{A - Correct:} Interdisciplinary team provides enterprise-wide view
				\item \textbf{B - Incorrect:} Without internal knowledge, assessment adequacy is difficult
				\item \textbf{C - Incorrect:} IT department unlikely to reflect entire enterprise
				\item \textbf{D - Incorrect:} Compliance department focuses on regulatory implications
			\end{itemize}
		\end{block}
	\end{frame}
	
	
	
	
	\begin{frame}{Q5: Enterprise Risk View} \begin{block}{Question} Executive management wants to obtain the most comprehensive view of enterprise risk before approving a major technology transformation initiative. Which approach would BEST support this objective? \end{block} \begin{enumerate}[(A)] \item Conduct a risk assessment using representatives from business, IT, legal, compliance, and operations \item Rely on an independent third-party risk assessment \item Obtain an IT risk assessment from the technology department \item Request a review from the compliance function \end{enumerate}  \begin{block}{CRISC Memory Aid} Enterprise Risk = Multiple Perspectives IT View = Technology Risk Compliance View = Regulatory Risk Third Party = Independent View Cross-Functional Team = Enterprise View \vspace{0.15cm} Exam Rule: When CRISC asks for the BEST enterprise-wide view, choose the answer involving multiple stakeholders. \end{block}  \end{frame} \begin{frame}{Q5: Explanation} \footnotesize \begin{block}{Correct Answer: A} Effective enterprise risk management requires input from multiple stakeholders. An interdisciplinary team provides perspectives on strategic, operational, legal, regulatory, financial, and technology risks, creating the most complete view of risk exposure. \end{block} \begin{block}{Option Analysis} \textbf{A - Correct} Enterprise risks often cross organizational boundaries. Business, IT, legal, compliance, operations, and other stakeholders each understand different risk exposures. Combining these viewpoints provides the most accurate assessment of enterprise risk. \vspace{0.15cm} \textbf{B - Incorrect} Third-party assessments can be valuable and provide independent insight. However, external assessors may lack detailed knowledge of internal processes, culture, strategy, and business objectives. \vspace{0.15cm} \textbf{C - Incorrect} IT assessments focus primarily on technology, cybersecurity, infrastructure, and application risks. Enterprise risk extends beyond technology and includes business, financial, operational, and regulatory concerns. \vspace{0.15cm} \textbf{D - Incorrect} Compliance reviews focus on adherence to laws, regulations, policies, and standards. While important, compliance represents only one dimension of enterprise risk. \end{block}\end{frame}
	
	
	\begin{frame}{Q6: IT Risk Awareness Program}
		\begin{block}{Question}
			Which is PRIMARY consideration when developing an IT risk awareness program?
		\end{block}
		\begin{enumerate}[(A)]
			\item Why technology risk is owned by IT
			\item How technology risk can impact each attendee's area of business
			\item How business process owners can transfer technology risk
			\item Why technology risk is more difficult to manage
		\end{enumerate}
	\end{frame}
	
	\begin{frame}{Q6: Explanation}
		\begin{block}{Correct Answer: B}
			Stakeholders must understand how the IT-related risk impacts the overall business.
		\end{block}
		\begin{block}{Option Analysis}
			\begin{itemize}
				\item \textbf{A - Incorrect:} IT does not own technology risk; business owns it
				\item \textbf{B - Correct:} Business impact is the primary consideration
				\item \textbf{C - Incorrect:} Transferring risk is part of risk response, not awareness
				\item \textbf{D - Incorrect:} Difficulty of management is secondary to business impact
			\end{itemize}
		\end{block}
	\end{frame}
	
	\begin{frame}{Q7: Risk Appetite Alignment}
		\begin{block}{Question}
			It is MOST important that risk appetite is aligned with business objectives to ensure that:
		\end{block}
		\begin{enumerate}[(A)]
			\item Resources are directed toward areas of low risk tolerance
			\item Major risk is identified and eliminated
			\item IT and business goals are aligned
			\item Risk strategy is adequately communicated
		\end{enumerate}
	\end{frame}
	
	\begin{frame}{Q7: Explanation}
		\begin{block}{Correct Answer: A}
			Risk appetite is the amount of risk that an enterprise is willing to take on in pursuit of value. Aligning it with business objectives allows an enterprise to evaluate and deploy valuable resources toward those objectives where the risk tolerance (for loss) is low.
		\end{block}
		\begin{block}{Option Analysis}
			\begin{itemize}
				\item \textbf{A - Correct:} Resources directed to low tolerance areas
				\item \textbf{B - Incorrect:} Risk elimination is rarely cost-effective
				\item \textbf{C - Incorrect:} Alignment affects more than IT and business
				\item \textbf{D - Incorrect:} Communication does not depend on alignment
			\end{itemize}
		\end{block}
	\end{frame}
	
	\begin{frame}[shrink=10]{Q7: Risk Appetite Alignment}
		\begin{block}{Question}
			Executive management is reviewing the organization's risk appetite statement to ensure it supports strategic growth initiatives. Why is it MOST important that risk appetite be aligned with business objectives?
		\end{block}
		\begin{enumerate}
			\item Resources can be prioritized toward objectives with the lowest tolerance for risk
			\item Significant risks can be completely eliminated
			\item Business and technology strategies remain synchronized
			\item Risk management expectations can be communicated consistently
		\end{enumerate}
		\begin{block}{CRISC Memory Aid}
			Risk Appetite = How Much Risk We Want\\
			Risk Tolerance = Allowed Variation\\
			Risk Capacity = Maximum Survivable Risk\\
			Strategy Drives Appetite\\
			Appetite Drives Investment\\
			Investment Drives Controls
		\end{block}
	\end{frame}
	
	\begin{frame}[shrink=10]{Q7: Explanation}
		\footnotesize
		\begin{block}{Correct Answer: A}
			Risk appetite defines the amount of risk the organization is willing to accept while pursuing business objectives. Alignment ensures management can allocate resources appropriately and focus attention on objectives where risk exposure exceeds acceptable levels.
		\end{block}
		\begin{block}{Option Analysis}
			\textbf{A -- Correct} Risk appetite helps management determine where resources, controls, and oversight should be concentrated. Objectives with lower tolerance for risk typically require stronger protection and greater investment.
			

			\textbf{B -- Incorrect} Risk management seeks to manage risk within acceptable levels, not eliminate all risk. Eliminating every significant risk would often be impractical and inconsistent with achieving business value.
			

			\textbf{C -- Incorrect} Alignment between business and technology goals is important, but this is only one benefit of risk appetite. The concept applies across the entire enterprise, not only IT.
			

			\textbf{D -- Incorrect} Communication is important, but communication alone does not explain why risk appetite must align with business objectives. The primary purpose is to support decision-making and resource allocation.
		\end{block}
	\end{frame}

	
	\begin{frame}{Q8: Weak Passwords and Unprotected Transmission}
		\begin{block}{Question}
			Weak passwords and transmission over unprotected communication lines are examples of:
		\end{block}
		\begin{enumerate}[(A)]
			\item Vulnerabilities
			\item Threats
			\item Probabilities
			\item Impacts
		\end{enumerate}
	\end{frame}
	
	\begin{frame}{Q8: Explanation}
		\begin{block}{Correct Answer: A}
			Vulnerabilities represent characteristics of information resources that may be exploited by a threat.
		\end{block}
		\begin{block}{Option Analysis}
			\begin{itemize}
				\item \textbf{A - Correct:} Weak passwords and unprotected transmission are vulnerabilities
				\item \textbf{B - Incorrect:} Threats are circumstances with potential to cause harm
				\item \textbf{C - Incorrect:} Probabilities represent likelihood of occurrence
				\item \textbf{D - Incorrect:} Impacts represent outcome of threat exploiting vulnerability
			\end{itemize}
		\end{block}
	\end{frame}
	
	% ================ SECTION 2: DOMAIN 2 - IT RISK ASSESSMENT ================
	
	\section{Domain 2: IT Risk Assessment}
	
	\begin{frame}{Q9: Quantitative Risk Analysis Drawback}
		\begin{block}{Question}
			MOST significant drawback of using quantitative risk analysis instead of qualitative risk analysis is:
		\end{block}
		\begin{enumerate}[(A)]
			\item Lower objectivity
			\item Greater reliance on expertise
			\item Less management buy-in
			\item Higher cost
		\end{enumerate}
	\end{frame}
	
	\begin{frame}{Q9: Explanation}
		\begin{block}{Correct Answer: D}
			Quantitative risk analysis is generally more complex and, therefore, more costly than qualitative risk analysis.
		\end{block}
		\begin{block}{Option Analysis}
			\begin{itemize}
				\item \textbf{A - Incorrect:} Neither method is fully objective
				\item \textbf{B - Incorrect:} Both require personnel with good business understanding
				\item \textbf{C - Incorrect:} Quantitative generally has better buy-in
				\item \textbf{D - Correct:} Quantitative is more complex and costly
			\end{itemize}
		\end{block}
	\end{frame}
	
	\begin{frame}{Q9: Quantitative Risk Analysis} \begin{block}{Question} An organization wants to perform a quantitative risk assessment for a new online banking platform. The risk practitioner notes that loss estimates, threat frequency data, and historical incident metrics must be collected before analysis can begin. What is the MOST significant challenge associated with using quantitative risk analysis? \end{block} \begin{enumerate}[(A)] \item Lower objectivity \item Greater reliance on expertise \item Reduced management support \item Higher cost and effort \end{enumerate}  \begin{block}{CRISC Memory Aid} Qualitative = Fast + Cheap + Subjective Quantitative = Slow + Expensive + Numerical \vspace{0.15cm} Qualitative: High / Medium / Low Quantitative: \$ Value, ALE, SLE, ARO \end{block} \end{frame} \begin{frame}{Q9: Explanation} \footnotesize \begin{block}{Correct Answer: D - Higher Cost and Effort} Quantitative risk analysis requires reliable data, statistical techniques, and significant analysis effort. Collecting and validating this information can be expensive and time-consuming, making quantitative approaches costlier than qualitative methods. \end{block} \begin{block}{Option Analysis} \textbf{A - Incorrect} Quantitative analysis is generally considered more objective because it relies on numerical estimates and measurable criteria. Qualitative analysis is typically more dependent on judgment. \vspace{0.15cm} \textbf{B - Incorrect} Both qualitative and quantitative approaches require experienced personnel. The difference is that quantitative analysis requires more detailed data collection and modeling. \vspace{0.15cm} \textbf{C - Incorrect} Management often prefers quantitative results because financial impacts and loss estimates can be expressed in business terms. \vspace{0.15cm} \textbf{D - Correct} Quantitative methods require historical loss information, probability estimates, asset valuations, and detailed analysis. Obtaining and maintaining this information increases cost, complexity, and effort. \end{block} \end{frame}
	
	
	\begin{frame}{Q10: Risk Scenario Analysis Purpose}
		\begin{block}{Question}
			Risk scenarios are analyzed to determine:
		\end{block}
		\begin{enumerate}[(A)]
			\item Strength of controls
			\item Likelihood and impact
			\item Current risk profile
			\item Scenario root cause
		\end{enumerate}
	\end{frame}
	
	\begin{frame}{Q10: Explanation}
		\begin{block}{Correct Answer: B}
			Risk scenarios are descriptions of events that can lead to a business impact and are evaluated to determine the likelihood and impact should the event occur.
		\end{block}
		\begin{block}{Option Analysis}
			\begin{itemize}
				\item \textbf{A - Incorrect:} Control strength is determined after controls are in place
				\item \textbf{B - Correct:} Likelihood and impact are the key assessments
				\item \textbf{C - Incorrect:} Risk profile is identification of current risk concerns
				\item \textbf{D - Incorrect:} Root cause is not part of risk scenario process
			\end{itemize}
		\end{block}
	\end{frame}
	
	\begin{frame}{Q11: Risk Ownership of Critical System}
		\begin{block}{Question}
			The risk to an information system that supports a critical business process is owned by:
		\end{block}
		\begin{enumerate}[(A)]
			\item IT director
			\item Senior management
			\item Risk management department
			\item System users
		\end{enumerate}
	\end{frame}
	
	\begin{frame}{Q11: Explanation}
		\begin{block}{Correct Answer: B}
			Senior management is responsible for the acceptance and mitigation of all risk.
		\end{block}
		\begin{block}{Option Analysis}
			\begin{itemize}
				\item \textbf{A - Incorrect:} IT director manages systems on behalf of business owners
				\item \textbf{B - Correct:} Senior management owns risk
				\item \textbf{C - Incorrect:} Risk management determines and reports on risk but does not own it
				\item \textbf{D - Incorrect:} System users use system properly but do not own risk
			\end{itemize}
		\end{block}
	\end{frame}
	
	
	
	
	\begin{frame}[shrink=10]{Q11: Risk Ownership}
		\begin{block}{Question}
			A critical payment-processing application supports the
			organization's primary revenue-generating business process.
			A significant risk associated with the application remains
			after controls have been implemented.
			
			Who should ultimately decide whether to accept this residual risk?
		\end{block}
		
		\begin{enumerate}[(A)]
			\item IT director
			\item Senior management
			\item Risk management department
			\item System administrator
		\end{enumerate}
		
				\begin{block}{CRISC Memory Aid}
			
			Business Owns Risk
			
			IT Owns Technology
			
			Security Owns Controls
			
			Risk Management Owns Process
			
			Audit Owns Assurance
			
			\vspace{0.15cm}
			
			Exam Rule:
			
			Responsibility can be delegated.
			
			Accountability cannot be delegated.
			
			Risk Acceptance = Management.
			
		\end{block}
	\end{frame}
	
	\begin{frame}[shrink=10]{Q11: Explanation}
		\footnotesize
		
		\begin{block}{Correct Answer: B - Senior Management}
			Management owns business risk and is ultimately responsible
			for determining whether residual risk remains within the
			organization's risk appetite and tolerance levels.
		\end{block}
		
		\begin{block}{Option Analysis}
			
			\textbf{A - Incorrect: IT Director}
			
			The IT director may implement controls, manage technology,
			and provide risk information. However, IT does not own the
			business risk created by the process.
			
			\vspace{0.15cm}
			
			\textbf{B - Correct: Senior Management}
			
			Senior management is accountable for risk decisions and risk
			acceptance. Because risk affects business objectives,
			management must determine whether residual risk is acceptable.
			
			\vspace{0.15cm}
			
			\textbf{C - Incorrect: Risk Management Department}
			
			Risk management facilitates assessments, monitoring,
			reporting, and governance activities. The department advises
			management but does not own business risk.
			
			\vspace{0.15cm}
			
			\textbf{D - Incorrect: System Administrator}
			
			System administrators operate and maintain systems. They may
			implement controls but do not have authority to accept
			enterprise risk on behalf of the organization.
			
		\end{block}
		

		
	\end{frame}
	
	\begin{frame}{Q12: Periodic Risk Assessment Reason}
		\begin{block}{Question}
			PRIMARY reason risk assessments should be repeated at regular intervals is:
		\end{block}
		\begin{enumerate}[(A)]
			\item Omissions in earlier assessments can be addressed
			\item Periodic assessments allow various methodologies
			\item Business threats are constantly changing
			\item They help raise risk awareness among staff
		\end{enumerate}
	\end{frame}
	
	\begin{frame}{Q12: Explanation}
		\begin{block}{Correct Answer: C}
			As business objectives and methods change, the nature and relevance of threats also change. This is the primary reason to conduct periodic risk assessments.
		\end{block}
		\begin{block}{Option Analysis}
			\begin{itemize}
				\item \textbf{A - Incorrect:} This is a benefit but not the primary reason
				\item \textbf{B - Incorrect:} Organizations strive to improve but this is not primary
				\item \textbf{C - Correct:} Threats constantly change
				\item \textbf{D - Incorrect:} Risk awareness is a minor benefit
			\end{itemize}
		\end{block}
	\end{frame}
	
	\begin{frame}[shrink=10]{Q13: Measuring Risk Management Development}
		\begin{block}{Question}
			Which BEST assists in measuring existing level of development of risk management processes against desired state?
		\end{block}
		\begin{enumerate}[(A)]
			\item Capability maturity model (CMM)
			\item Risk management audit reports
			\item Balanced scorecard
			\item Enterprise security architecture
		\end{enumerate}
	\end{frame}
	
	\begin{frame}{Q13: Explanation}
		\begin{block}{Correct Answer: A}
			The capability maturity model (CMM) grades processes on a scale of 0 to 5, based on their maturity. It is commonly used by entities to measure their existing state and then to determine the desired one.
		\end{block}
		\begin{block}{Option Analysis}
			\begin{itemize}
				\item \textbf{A - Correct:} CMM measures current vs desired state
				\item \textbf{B - Incorrect:} Audit reports offer limited view
				\item \textbf{C - Incorrect:} Balanced scorecard measures implementation of strategy
				\item \textbf{D - Incorrect:} Enterprise security architecture explains business strategy view
			\end{itemize}
		\end{block}
	\end{frame}
	
	
	
	
	
	
	
	
	
	
	\begin{frame}{Q13: Risk Management Maturity} \begin{block}{Question} A Chief Risk Officer wants to determine how mature the organization's risk management program is and compare the current state against a targeted future state approved by executive management. Which would BEST support this assessment? \end{block} \begin{enumerate}[(A)] \item Capability Maturity Model (CMM) \item Risk management audit reports \item Balanced scorecard \item Enterprise security architecture \end{enumerate}  \begin{block}{CRISC Memory Aid} CMM = "How mature are we?" Audit = "What is wrong?" Balanced Scorecard = "Are we meeting strategy?" Security Architecture = "How is security designed?" \vspace{0.15cm} CMM Levels: Level 1 = Initial (Ad Hoc) Level 2 = Repeatable Level 3 = Defined Level 4 = Managed (Measured) Level 5 = Optimized (Continuous Improvement) \end{block} \end{frame} \begin{frame}{Q13: Explanation} \footnotesize \begin{block}{Correct Answer: A - Capability Maturity Model (CMM)} Maturity models provide a structured method for evaluating current process capability and comparing it against a desired future state. This enables organizations to identify gaps and prioritize improvement initiatives. \end{block} \begin{block}{Option Analysis} \textbf{A - Correct} CMM measures process maturity across progressive levels. Organizations commonly use maturity assessments to identify where they are today, where they want to be, and what improvements are needed to reach the desired state. \vspace{0.15cm} \textbf{B - Incorrect} Audit reports identify control deficiencies and compliance issues. While useful for improvement, they do not provide a comprehensive maturity benchmark or roadmap. \vspace{0.15cm} \textbf{C - Incorrect} A balanced scorecard measures organizational performance against strategic objectives. It is not primarily designed to evaluate process maturity. \vspace{0.15cm} \textbf{D - Incorrect} Enterprise security architecture describes how security capabilities, systems, and controls fit together. It does not measure the maturity of risk management processes. \end{block} \end{frame}
	
	
	
	
	\begin{frame}{Q14: Identifying Control Deficiencies}
		\begin{block}{Question}
			Which BEST helps identify information systems control deficiencies?
		\end{block}
		\begin{enumerate}[(A)]
			\item Gap analysis
			\item Current IT risk profile
			\item IT controls framework
			\item Countermeasure analysis
		\end{enumerate}
	\end{frame}
	
	\begin{frame}{Q14: Explanation}
		\begin{block}{Correct Answer: A}
			Controls are deployed to achieve desired control objectives based on risk assessments and business requirements. The gap between desired control objectives and actual IS control design and operational effectiveness identifies control deficiencies.
		\end{block}
		\begin{block}{Option Analysis}
			\begin{itemize}
				\item \textbf{A - Correct:} Gap analysis identifies deficiencies
				\item \textbf{B - Incorrect:} Risk profile does not expose gaps
				\item \textbf{C - Incorrect:} Framework is generic without current state info
				\item \textbf{D - Incorrect:} Only identifies countermeasure deficiencies
			\end{itemize}
		\end{block}
	\end{frame}
	
	
	
	
	
	
	
	
	
	\begin{frame}{Q14: Identifying Control Deficiencies}
		\begin{block}{Question}
			A risk practitioner is reviewing access management controls for a
			critical financial application. Management wants to determine
			whether existing controls adequately satisfy policy requirements
			and regulatory expectations.
			
			Which technique would BEST identify control deficiencies?
		\end{block}
		
		\begin{enumerate}[(A)]
			\item Gap analysis
			\item Current IT risk profile
			\item IT controls framework
			\item Countermeasure analysis
		\end{enumerate}
		
			\begin{block}{CRISC Memory Aid}
			
			Gap Analysis
			
			= Current State vs Desired State
			
			= Missing Controls
			
			= Weak Controls
			
			= Process Deficiencies
			
			\vspace{0.15cm}
			
			CMM
			
			= Maturity Gap
			
			Gap Analysis
			
			= Control Gap
			
			Risk Assessment
			
			= Risk Gap
			
		\end{block}
	\end{frame}
	
	\begin{frame}[shrink=10]{Q14: Explanation}
		\footnotesize
		
		\begin{block}{Correct Answer: A - Gap Analysis}
			A gap analysis compares the current control environment against
			desired control objectives, regulatory requirements, policies,
			or industry standards. Any difference between the current state
			and the desired state represents a potential control deficiency.
		\end{block}
		
		\begin{block}{Option Analysis}
			
			\textbf{A - Correct}
			
			Gap analysis is specifically designed to identify missing,
			inadequate, or ineffective controls. It helps determine what
			must be improved to meet organizational requirements.
			
			\vspace{0.15cm}
			
			\textbf{B - Incorrect}
			
			A current IT risk profile provides a snapshot of identified risks
			and exposures. While useful for understanding risk, it does not
			directly identify which controls are missing or ineffective.
			
			\vspace{0.15cm}
			
			\textbf{C - Incorrect}
			
			An IT controls framework such as COBIT or NIST provides guidance
			on what controls should exist. However, the framework alone does
			not reveal whether an organization's current controls meet those
			expectations.
			
			\vspace{0.15cm}
			
			\textbf{D - Incorrect}
			
			Countermeasure analysis focuses on evaluating potential safeguards
			after risks or control gaps have already been identified. It is
			not the primary technique for finding deficiencies.
			
		\end{block}
		
	
		
	\end{frame}
	
	
	\begin{frame}{Q15: Statistical Methods Risk Analysis}
		\begin{block}{Question}
			Deriving likelihood and impact of risk scenarios through statistical methods is MOST LIKELY associated with which type of risk analysis?
		\end{block}
		\begin{enumerate}[(A)]
			\item Risk scenario
			\item Qualitative
			\item Quantitative
			\item Semiquantitative
		\end{enumerate}
	\end{frame}
	
	\begin{frame}{Q15: Explanation}
		\begin{block}{Correct Answer: C}
			The essence of quantitative risk assessment is to derive the likelihood and impact of risk scenarios based on statistical methods and data.
		\end{block}
		\begin{block}{Option Analysis}
			\begin{itemize}
				\item \textbf{A - Incorrect:} Scenario analysis includes several methods
				\item \textbf{B - Incorrect:} Qualitative uses experiential measures
				\item \textbf{C - Correct:} Quantitative uses statistical methods
				\item \textbf{D - Incorrect:} Semiquantitative assigns values to qualitative ratings
			\end{itemize}
		\end{block}
	\end{frame}
	
	\begin{frame}[shrink=10]{Q16: Reviewing Risk Analysis Results}
		\begin{block}{Question}
			Which review is BEST suited for review of IT risk analysis results before sending to management?
		\end{block}
		\begin{enumerate}[(A)]
			\item Internal audit review
			\item Peer review
			\item Compliance review
			\item Risk policy review
		\end{enumerate}
	\end{frame}
	
	\begin{frame}[shrink=10]{Q16: Explanation}
		\begin{block}{Correct Answer: B}
			It is effective, efficient and good practice to perform a peer review of IT risk analysis results before sending them to management.
		\end{block}
		\begin{block}{Option Analysis}
			\begin{itemize}
				\item \textbf{A - Incorrect:} Audit is independent assurance
				\item \textbf{B - Correct:} Peer review is good practice
				\item \textbf{C - Incorrect:} Compliance reviews measure conformance
				\item \textbf{D - Incorrect:} Policy review may change content but isn't a review method
			\end{itemize}
		\end{block}
		
	\end{frame}
	
	
		\begin{frame}[shrink=10]{Q16: Reviewing Risk Analysis Results}
		\begin{block}{Question}
			Which review is BEST suited for review of IT risk analysis results before sending to management?
		\end{block}
		\begin{enumerate}[(A)]
			\item Internal audit review
			\item Peer review
			\item Compliance review
			\item Risk policy review
		\end{enumerate}
	\end{frame}
	
	\begin{frame}{Q16: Explanation}
		\begin{block}{Correct Answer: B}
			It is effective, efficient and good practice to perform a peer review of IT risk analysis results before sending them to management.
		\end{block}

	\end{frame}
	
	
	% ================ SECTION 3: DOMAIN 3 - RISK RESPONSE ================
	
	\section{Domain 3: Risk Response and Mitigation}
	
	\begin{frame}{Q17: Risk Cannot Be Sufficiently Mitigated}
		\begin{block}{Question}
			When risk cannot be sufficiently mitigated through manual or automatic controls, which option BEST protects from potential financial impact?
		\end{block}
		\begin{enumerate}[(A)]
			\item Insuring against the risk
			\item Updating IT risk register
			\item Improving staff training
			\item Outsourcing related process
		\end{enumerate}
	\end{frame}
	
	\begin{frame}{Q17: Explanation}
		\begin{block}{Correct Answer: A}
			An insurance policy can compensate the enterprise monetarily for the impact of the risk by transferring the risk to the insurance company.
		\end{block}
		\begin{block}{Option Analysis}
			\begin{itemize}
				\item \textbf{A - Correct:} Insurance transfers financial risk
				\item \textbf{B - Incorrect:} Updating register doesn't change risk
				\item \textbf{C - Incorrect:} Staff training may reduce but not completely mitigate
				\item \textbf{D - Incorrect:} Outsourcing does not necessarily remove risk
			\end{itemize}
		\end{block}
	\end{frame}
	
	\begin{frame}{Q18: Risk Mitigation Must Reduce}
		\begin{block}{Question}
			To be effective, risk mitigation MUST reduce the:
		\end{block}
		\begin{enumerate}[(A)]
			\item Residual risk
			\item Inherent risk
			\item Frequency of a threat
			\item Impact of a threat
		\end{enumerate}
	\end{frame}
	
	\begin{frame}{Q18: Explanation}
		\begin{block}{Correct Answer: A}
			The objective of risk reduction is to reduce the residual risk to levels below the enterprise's risk tolerance level.
		\end{block}
		\begin{block}{Option Analysis}
			\begin{itemize}
				\item \textbf{A - Correct:} Residual risk must be reduced to acceptable level
				\item \textbf{B - Incorrect:} Reduction of inherent risk is not the goal
				\item \textbf{C - Incorrect:} Can focus on reducing either frequency or impact
				\item \textbf{D - Incorrect:} Can focus on reducing either frequency or impact
			\end{itemize}
		\end{block}
	\end{frame}
	
	\begin{frame}{Q19: Best Control to Prevent Unauthorized Access}
		\begin{block}{Question}
			BEST control to prevent unauthorized access to information is user:
		\end{block}
		\begin{enumerate}[(A)]
			\item Accountability
			\item Authentication
			\item Identification
			\item Access rules
		\end{enumerate}
	\end{frame}
	
	\begin{frame}{Q19: Explanation}
		\begin{block}{Correct Answer: B}
			Authentication verifies the user's identity and the right to access information according to the access rules.
		\end{block}
		\begin{block}{Option Analysis}
			\begin{itemize}
				\item \textbf{A - Incorrect:} Accountability maps activity back to responsible party
				\item \textbf{B - Correct:} Authentication verifies identity
				\item \textbf{C - Incorrect:} Identification without authentication does not grant access
				\item \textbf{D - Incorrect:} Access rules without identification and authentication do not grant access
			\end{itemize}
		\end{block}
	\end{frame}
	
	
	
	\begin{frame}{Q19: Preventing Unauthorized Access} \begin{block}{Question} During a review of a critical payroll application, management discovers that several terminated employees attempted to log in using previously assigned user IDs. Which control is MOST important for preventing those individuals from gaining access to the system? \end{block} \begin{enumerate}[(A)] \item Accountability \item Authentication \item Identification \item Access rules \end{enumerate} \end{frame} \begin{frame}[shrink=10]{Q19: Explanation} \footnotesize \begin{block}{Correct Answer: B - Authentication} Authentication verifies that a user is who they claim to be before access is granted. Without successful authentication, unauthorized individuals cannot access protected resources. \end{block} \begin{block}{Option Analysis} \textbf{A - Incorrect: Accountability} Accountability ensures user actions can be traced to a specific individual through logs, audit trails, and monitoring. It helps investigate incidents but does not prevent access. \vspace{0.15cm} \textbf{B - Correct: Authentication} Authentication validates a user's identity through passwords, MFA, smart cards, biometrics, or certificates. It is the primary control used to prevent unauthorized users from accessing systems. \vspace{0.15cm} \textbf{C - Incorrect: Identification} Identification occurs when a user claims an identity by entering a user ID or account name. Anyone can claim an identity, so identification alone does not prevent unauthorized access. \vspace{0.15cm} \textbf{D - Incorrect: Access Rules} Access rules define what authenticated users are allowed to do. However, access rules are applied only after identity has been verified through authentication. \end{block}  \end{frame}
	
	\begin{frame}{Q20: Controlling Access to Sensitive Information}
		\begin{block}{Question}
			Which control BEST protects from unauthorized individuals gaining access to sensitive information?
		\end{block}
		\begin{enumerate}[(A)]
			\item Using a challenge response system
			\item Forcing periodic password changes
			\item Monitoring unsuccessful login attempts
			\item Providing access on a need-to-know basis
		\end{enumerate}
	\end{frame}
	
	\begin{frame}{Q20: Explanation}
		\begin{block}{Correct Answer: D}
			Physical or logical system access should be assigned on a need-to-know basis (legitimate business requirements) and in ways that incorporate least privilege and segregation of duties.
		\end{block}
		\begin{block}{Option Analysis}
			\begin{itemize}
				\item \textbf{A - Incorrect:} Challenge response doesn't address access design
				\item \textbf{B - Incorrect:} Password changes don't guarantee appropriate access
				\item \textbf{C - Incorrect:} Monitoring doesn't address appropriate access rights
				\item \textbf{D - Correct:} Need-to-know basis is fundamental
			\end{itemize}
		\end{block}
	\end{frame}
	
	\begin{frame}{Q21: Best Defense Against Phishing}
		\begin{block}{Question}
			Which defense is BEST to use against phishing attacks?
		\end{block}
		\begin{enumerate}[(A)]
			\item Intrusion detection system
			\item Spam filters
			\item End-user awareness
			\item Application hardening
		\end{enumerate}
	\end{frame}
	
	\begin{frame}{Q21: Explanation}
		\begin{block}{Correct Answer: C}
			Phishing attacks are a type of social engineering attack and are best defended by end-user awareness training.
		\end{block}
		\begin{block}{Option Analysis}
			\begin{itemize}
				\item \textbf{A - Incorrect:} IDS doesn't protect against phishing
				\item \textbf{B - Incorrect:} Spam filters are not as effective as user awareness
				\item \textbf{C - Correct:} User awareness is the best defense
				\item \textbf{D - Incorrect:} Application hardening doesn't protect against phishing
			\end{itemize}
		\end{block}
	\end{frame}
	
	\begin{frame}{Q22: Stakeholders in Risk Response Review}
		\begin{block}{Question}
			When responding to identified risk, MOST important stakeholders involved in reviewing risk response options are:
		\end{block}
		\begin{enumerate}[(A)]
			\item Information security managers
			\item Internal auditors
			\item Incident response team members
			\item Business managers
		\end{enumerate}
	\end{frame}
	
	\begin{frame}{Q22: Explanation}
		\begin{block}{Correct Answer: D}
			Business managers are accountable for managing the associated risk and will determine what actions to take based on the information provided by others, which may include collaboration with and support from IT security managers.
		\end{block}
		\begin{block}{Option Analysis}
			\begin{itemize}
				\item \textbf{A - Incorrect:} Security managers understand technical aspects but don't decide
				\item \textbf{B - Incorrect:} Risk response is not a function of internal audit
				\item \textbf{C - Incorrect:} Incident response team responds to incidents
				\item \textbf{D - Correct:} Business managers are accountable
			\end{itemize}
		\end{block}
	\end{frame}
	
	\begin{frame}{Q23: Designing Information System Controls}
		\begin{block}{Question}
			Which should be considered FIRST when designing information system controls?
		\end{block}
		\begin{enumerate}[(A)]
			\item Organizational strategic plan
			\item Existing IT environment
			\item Present IT budget
			\item IT strategic plan
		\end{enumerate}
	\end{frame}
	
	\begin{frame}{Q23: Explanation}
		\begin{block}{Correct Answer: A}
			Review of the enterprise's strategic plan is the first step in designing effective IS controls that would fit the enterprise's long-term plans.
		\end{block}
		\begin{block}{Option Analysis}
			\begin{itemize}
				\item \textbf{A - Correct:} Strategic plan is the starting point
				\item \textbf{B - Incorrect:} Useful but not the first task
				\item \textbf{C - Incorrect:} Budget is one component of strategic plan
				\item \textbf{D - Incorrect:} IT strategic plan exists to support enterprise's strategic plan
			\end{itemize}
		\end{block}
	\end{frame}
	
	\begin{frame}{Q24: Residual Risk Calculation Basis}
		\begin{block}{Question}
			Residual risk can be accurately calculated on the basis of:
		\end{block}
		\begin{enumerate}[(A)]
			\item Threats and vulnerabilities
			\item Inherent risk and control risk
			\item Compliance risk and reputation
			\item Risk governance and risk response
		\end{enumerate}
	\end{frame}
	
	\begin{frame}{Q24: Explanation}
		\begin{block}{Correct Answer: B}
			Inherent risk multiplied by control risk is the formula to calculate residual risk.
		\end{block}
		\begin{block}{Option Analysis}
			\begin{itemize}
				\item \textbf{A - Incorrect:} Threat/vulnerability are elements of inherent risk
				\item \textbf{B - Correct:} Inherent risk × control risk = residual risk
				\item \textbf{C - Incorrect:} Compliance risk is not a basis for calculation
				\item \textbf{D - Incorrect:} These are risk domains, not elements
			\end{itemize}
		\end{block}
	\end{frame}
	
	
	\begin{frame}{Q24: Residual Risk} \begin{block}{Question} A risk assessment determined that a critical application has a high inherent risk rating. After evaluating the effectiveness of existing controls, management wants to determine the remaining level of risk exposure. Which factors are MOST important when determining residual risk? \end{block} \begin{enumerate}[(A)] \item Threats and vulnerabilities \item Inherent risk and control risk \item Compliance risk and reputational risk \item Risk governance and risk response \end{enumerate} \begin{block}{CRISC Memory Aid} Inherent Risk = Risk Before Controls \vspace{0.1cm} Residual Risk = Risk After Controls \vspace{0.1cm} Strong Controls $\rightarrow$ Lower Residual Risk Weak Controls $\rightarrow$ Higher Residual Risk \vspace{0.15cm} Exam Rule: Threats + Vulnerabilities $\rightarrow$ Inherent Risk Inherent Risk + Control Effectiveness $\rightarrow$ Residual Risk \end{block} \end{frame} \begin{frame}[shrink=10]{Q24: Explanation} \footnotesize \begin{block}{Correct Answer: B - Inherent Risk and Control Risk} Residual risk represents the amount of risk remaining after controls have been implemented. It is determined by evaluating the inherent risk and the effectiveness of the controls used to reduce that risk. \end{block} \begin{block}{Option Analysis} \textbf{A - Incorrect} Threats and vulnerabilities contribute to inherent risk but do not by themselves determine how much risk remains after controls are applied. \vspace{0.15cm} \textbf{B - Correct} Residual risk is derived from the level of inherent risk and the effectiveness of the controls mitigating that risk. Weak controls result in higher residual risk, while effective controls reduce residual risk. \vspace{0.15cm} \textbf{C - Incorrect} Compliance and reputational risks are risk categories that may be affected by an event, but they are not the primary inputs used to determine residual risk. \vspace{0.15cm} \textbf{D - Incorrect} Governance and risk response activities influence how risks are managed but are not direct components of residual risk calculation. \end{block}  \end{frame}
	
	% ================ SECTION 4: DOMAIN 4 - RISK MONITORING ================
	
	\section{Domain 4: Risk and Control Monitoring}
	
	\begin{frame}{Q25: Maintaining KRIs}
		\begin{block}{Question}
			MOST important reason to maintain key risk indicators (KRIs) is that:
		\end{block}
		\begin{enumerate}[(A)]
			\item Complex metrics require fine-tuning
			\item Threats and vulnerabilities change over time
			\item Risk reports need to be timely
			\item They help to avoid risk
		\end{enumerate}
	\end{frame}
	
	\begin{frame}{Q25: Explanation}
		\begin{block}{Correct Answer: B}
			Threats and vulnerabilities change over time, and KRI maintenance ensures that KRIs continue to effectively capture these changes.
		\end{block}
		\begin{block}{Option Analysis}
			\begin{itemize}
				\item \textbf{A - Incorrect:} Fine-tuning is secondary to capturing changes
				\item \textbf{B - Correct:} Threats and vulnerabilities change over time
				\item \textbf{C - Incorrect:} Timeliness is a business requirement
				\item \textbf{D - Incorrect:} Risk avoidance is one risk response
			\end{itemize}
		\end{block}
	\end{frame}
	
	\begin{frame}{Q26: Measuring Operational Effectiveness}
		\begin{block}{Question}
			Which is BEST measure of operational effectiveness of risk management process capabilities?
		\end{block}
		\begin{enumerate}[(A)]
			\item Key performance indicators (KPIs)
			\item Key risk indicators (KRIs)
			\item Base practices
			\item Metric thresholds
		\end{enumerate}
	\end{frame}
	
	\begin{frame}{Q26: Explanation}
		\begin{block}{Correct Answer: A}
			Key performance indicators (KPIs) are assessment indicators that support judgment regarding the performance of a specific process.
		\end{block}
		\begin{block}{Option Analysis}
			\begin{itemize}
				\item \textbf{A - Correct:} KPIs measure process performance
				\item \textbf{B - Incorrect:} KRIs provide insights into potential risk
				\item \textbf{C - Incorrect:} Base practices need work products for evidence
				\item \textbf{D - Incorrect:} Thresholds are decision points
			\end{itemize}
		\end{block}
	\end{frame}
	
	\begin{frame}{Q27: Data Extraction Forecasting}
		\begin{block}{Question}
			During data extraction, total transactions per year was forecasted by multiplying monthly average by twelve. This is considered:
		\end{block}
		\begin{enumerate}[(A)]
			\item A controls total
			\item Simplistic and ineffective
			\item A duplicates test
			\item A reasonableness test
		\end{enumerate}
	\end{frame}
	
	\begin{frame}{Q27: Explanation}
		\begin{block}{Correct Answer: D}
			Reasonableness tests make certain assumptions about the information as the basis for more elaborate data validation tests.
		\end{block}
		\begin{block}{Option Analysis}
			\begin{itemize}
				\item \textbf{A - Incorrect:} Does not ensure all transactions extracted
				\item \textbf{B - Incorrect:} Reasonableness test is valid foundation
				\item \textbf{C - Incorrect:} Does not identify duplicate transactions
				\item \textbf{D - Correct:} This is a reasonableness test
			\end{itemize}
		\end{block}
	\end{frame}
	
	
	\begin{frame}{Q27: Validating Extracted Data}
		\begin{block}{Question}
			A risk practitioner extracts one year of transaction data from a
			claims-processing system for analysis. To verify the completeness
			of the extract, the practitioner compares the extracted annual
			transaction count to an estimate derived by multiplying the
			average monthly transaction volume by twelve.
			
			What type of test is being performed?
		\end{block}
		
		\begin{enumerate}[(A)]
			\item Control total test
			\item Ineffective forecasting technique
			\item Duplicate transaction test
			\item Reasonableness test
		\end{enumerate}
		
			
		\begin{block}{CRISC Memory Aid}
			
			Control Total
			
			= Did we get all the data?
			
			\vspace{0.1cm}
			
			Duplicate Test
			
			= Did we get the same data twice?
			
			\vspace{0.1cm}
			
			Reasonableness Test
			
			= Does the result make sense?
			
			\vspace{0.15cm}
			
			If you see:
			
			"Expected vs Actual"
			
			think
			
			"Reasonableness Test"
			
		\end{block}
		
	\end{frame}
	
	\begin{frame}[shrink=10]{Q27: Explanation}
		\footnotesize
		
		\begin{block}{Correct Answer: D - Reasonableness Test}
			A reasonableness test compares actual results to an expected
			value derived from logical assumptions or historical trends.
			The purpose is to identify anomalies that may indicate data
			quality or extraction issues.
		\end{block}
		
		\begin{block}{Option Analysis}
			
			\textbf{A - Incorrect}
			
			Control totals verify completeness by comparing totals such as
			record counts, dollar amounts, or checksums between the source
			data and extracted data. The scenario uses an estimate rather
			than a control total.
			
			\vspace{0.15cm}
			
			\textbf{B - Incorrect}
			
			Although the estimate is simple, a reasonableness test does not
			require complex statistical methods. High-level estimates are
			commonly used to identify obvious discrepancies.
			
			\vspace{0.15cm}
			
			\textbf{C - Incorrect}
			
			Duplicate testing identifies repeated records or transactions.
			The scenario focuses on validating expected transaction volume,
			not identifying duplicate data.
			
			\vspace{0.15cm}
			
			\textbf{D - Correct}
			
			The practitioner is comparing actual results to a logical
			expectation based on historical averages. This is a classic
			reasonableness test.
			
		\end{block}
	
		
	\end{frame}
	
	
	\begin{frame}{Q28: Testing Access Management Effectiveness}
		\begin{block}{Question}
			BEST test for confirming effectiveness of system access management process is to map:
		\end{block}
		\begin{enumerate}[(A)]
			\item Access requests to user accounts
			\item User accounts to access requests
			\item User accounts to HR records
			\item Vendor database to user accounts
		\end{enumerate}
	\end{frame}
	
	\begin{frame}{Q28: Explanation}
		\begin{block}{Correct Answer: B}
			Mapping user accounts to access requests confirms that all existing accounts have been approved.
		\end{block}
		\begin{block}{Option Analysis}
			\begin{itemize}
				\item \textbf{A - Incorrect:} Confirms requests processed but not unapproved accounts
				\item \textbf{B - Correct:} Confirms all accounts have been approved
				\item \textbf{C - Incorrect:} Confirms accounts tied to employees only
				\item \textbf{D - Incorrect:} Flawed because it doesn't consider unapproved accounts
			\end{itemize}
		\end{block}
	\end{frame}
	
	\begin{frame}{Q29: Firewall Compliance Assurance}
		\begin{block}{Question}
			Which provides BEST assurance that a firewall is configured in compliance with enterprise's security policy?
		\end{block}
		\begin{enumerate}[(A)]
			\item Reviewing actual procedures
			\item Interviewing firewall administrator
			\item Reviewing parameter settings
			\item Reviewing log file for recent attacks
		\end{enumerate}
	\end{frame}
	
	\begin{frame}{Q29: Explanation}
		\begin{block}{Correct Answer: C}
			A review of the parameter settings provides a good basis for comparison of the actual configuration to the security policy and reliable audit evidence documentation.
		\end{block}
		\begin{block}{Option Analysis}
			\begin{itemize}
				\item \textbf{A - Incorrect:} Procedures show how it should be managed, not compliance
				\item \textbf{B - Incorrect:} Interview provides overview, not confirmation
				\item \textbf{C - Correct:} Parameter settings show actual configuration
				\item \textbf{D - Incorrect:} Log review shows logging is enabled, not configuration
			\end{itemize}
		\end{block}
	\end{frame}
	
	\begin{frame}{Q30: Verifying Control Effectiveness}
		\begin{block}{Question}
			One way to verify control effectiveness is by determining:
		\end{block}
		\begin{enumerate}[(A)]
			\item Its reliability
			\item Whether it is preventive or detective
			\item Capability of providing notification of failure
			\item Test results of intended objectives
		\end{enumerate}
	\end{frame}
	
	\begin{frame}{Q30: Explanation}
		\begin{block}{Correct Answer: D}
			Control effectiveness requires a process to verify that the control process worked as intended and meets the intended control objectives.
		\end{block}
		\begin{block}{Option Analysis}
			\begin{itemize}
				\item \textbf{A - Incorrect:} Weak controls can be highly reliable
				\item \textbf{B - Incorrect:} Control type doesn't determine effectiveness
				\item \textbf{C - Incorrect:} Notification of failure doesn't determine strength
				\item \textbf{D - Correct:} Test results verify intended objectives
			\end{itemize}
		\end{block}
	\end{frame}
	
	
	\begin{frame}{Q30: Verifying Control Effectiveness} \begin{block}{Question} A risk practitioner is evaluating a recently implemented access control to determine whether it is effectively reducing the targeted risk. Which provides the BEST evidence of control effectiveness? \end{block} \begin{enumerate}[(A)] \item Consistent operation of the control over time \item Testing that demonstrates achievement of control objectives \item Automated notification when the control fails \item Formal documentation of the control design \end{enumerate} \end{frame} \begin{frame}[shrink=10]{Q30: Explanation} \footnotesize \begin{block}{Correct Answer: B} Control effectiveness is determined by whether the control achieves its intended objective. Testing provides objective evidence that the control is mitigating risk as designed. \end{block} \begin{block}{Option Analysis} \textbf{A - Incorrect} Consistent operation indicates reliability, but a control can operate reliably and still fail to achieve its objective. \vspace{0.15cm} \textbf{B - Correct} Testing against control objectives provides direct evidence that the control is effective in reducing risk. \vspace{0.15cm} \textbf{C - Incorrect} Failure notifications support monitoring but do not prove that the control is effectively mitigating risk. \vspace{0.15cm} \textbf{D - Incorrect} Good design documentation indicates the control was planned appropriately, but documentation alone does not demonstrate operational effectiveness. \end{block} \end{frame}
	
	\begin{frame}{Q31: Correlation Tools for Trend Analysis}
		\begin{block}{Question}
			Tools that correlate information from multiple systems to improve trend analysis are MOST likely applied to:
		\end{block}
		\begin{enumerate}[(A)]
			\item Transaction data
			\item Configuration settings
			\item System changes
			\item Process integrity
		\end{enumerate}
	\end{frame}
	
	\begin{frame}{Q31: Explanation}
		\begin{block}{Correct Answer: A}
			Transaction data tends to be difficult to analyze for trends across multiple systems because of sheer volume. Therefore, correlation engines are often used to analyze the data to see trends that otherwise might not be apparent.
		\end{block}
		\begin{block}{Option Analysis}
			\begin{itemize}
				\item \textbf{A - Correct:} Transaction data benefits from correlation
				\item \textbf{B - Incorrect:} Configuration settings are compared to predefined values
				\item \textbf{C - Incorrect:} System changes rarely benefit from trend analysis
				\item \textbf{D - Incorrect:} Process integrity not typically associated with trend analysis
			\end{itemize}
		\end{block}
	\end{frame}
	
	
	\begin{frame}{Q31: Correlation Analysis} \begin{block}{Question} A financial institution implements a monitoring platform that collects events from payment systems, databases, web servers, and authentication services. The platform correlates information from these sources to identify unusual patterns that may indicate fraud or security incidents. This capability is MOST valuable when analyzing: \end{block} \begin{enumerate}[(A)] \item Transaction data \item Configuration settings \item System changes \item Process integrity \end{enumerate} \end{frame} \begin{frame}[shrink=10]{Q31: Explanation} \footnotesize \begin{block}{Correct Answer: A - Transaction Data} Correlation engines are most effective when processing large volumes of transactional information from multiple systems. By linking related events across platforms, organizations can identify trends, anomalies, fraud indicators, and emerging risks that may not be visible when each system is analyzed separately. \end{block} \begin{block}{Option Analysis} \textbf{A - Correct} Transaction processing generates large volumes of data across multiple systems. Correlation tools can connect seemingly unrelated events to identify suspicious patterns, operational issues, or emerging risk trends. \vspace{0.15cm} \textbf{B - Incorrect} Configuration settings are typically compared against approved baselines, standards, or secure configurations. Correlation analysis provides less value than configuration-compliance reviews. \vspace{0.15cm} \textbf{C - Incorrect} System changes are generally reviewed through change-management processes, approvals, and audit trails. While correlation may provide supporting information, it is not the primary use case. \vspace{0.15cm} \textbf{D - Incorrect} Process integrity is usually assessed through controls, monitoring, segregation of duties, and audits rather than large-scale event correlation. \end{block} \begin{block}{CRISC Memory Aid} SIEM = Correlate Events = Find Patterns = Detect Trends \vspace{0.15cm} Large Volume Data $\rightarrow$ Correlation Transaction Monitoring $\rightarrow$ Correlation Fraud Analytics $\rightarrow$ Correlation \end{block} \end{frame}
	
	\begin{frame}{Q32: Ensuring Outsourced Compliance}
		\begin{block}{Question}
			Most effective way to ensure outsourced service providers comply with enterprise information security policy?
		\end{block}
		\begin{enumerate}[(A)]
			\item Periodic audits
			\item Security awareness training
			\item Penetration testing
			\item Service level monitoring
		\end{enumerate}
	\end{frame}
	
	\begin{frame}{Q32: Explanation}
		\begin{block}{Correct Answer: A}
			Regular audits can identify gaps in information security compliance.
		\end{block}
		\begin{block}{Option Analysis}
			\begin{itemize}
				\item \textbf{A - Correct:} Audits ensure compliance
				\item \textbf{B - Incorrect:} Training increases awareness but less effective
				\item \textbf{C - Incorrect:} Penetration testing identifies vulnerabilities
				\item \textbf{D - Incorrect:} Service monitoring identifies operational issues
			\end{itemize}
		\end{block}
	\end{frame}
	
	
	
	\begin{frame}{Q32: Third-Party Security Compliance} \begin{block}{Question} An organization outsources payroll processing to a third-party service provider. Management wants ongoing assurance that the provider is complying with the organization's information security requirements. Which approach would provide the MOST effective assurance? \end{block} \begin{enumerate}[(A)] \item Periodic independent security audits \item Annual security awareness training for vendor personnel \item Regular penetration testing of the provider environment \item Monitoring service level agreement (SLA) metrics \end{enumerate} \end{frame} \begin{frame}[shrink=10]{Q32: Explanation} \footnotesize \begin{block}{Correct Answer: A - Periodic Independent Security Audits} Independent audits provide objective evidence regarding whether the service provider is complying with contractual, regulatory, and security requirements. Audits evaluate both the design and operating effectiveness of controls. \end{block} \begin{block}{Option Analysis} \textbf{A - Correct} Audits provide formal assurance that controls are operating as intended and that security requirements are being met. This is the strongest mechanism for verifying compliance. \vspace{0.15cm} \textbf{B - Incorrect} Security awareness training improves knowledge and behavior, but training alone does not verify compliance with security requirements. \vspace{0.15cm} \textbf{C - Incorrect} Penetration testing identifies technical vulnerabilities and security weaknesses. A provider may pass a penetration test while still failing to comply with policy requirements. \vspace{0.15cm} \textbf{D - Incorrect} SLA monitoring evaluates operational performance such as uptime, response times, and service quality. It does not directly assess security-policy compliance. \end{block} \begin{block}{CRISC Memory Aid} Audit = Compliance Assurance \vspace{0.1cm} Pen Test = Technical Weaknesses \vspace{0.1cm} Training = Awareness \vspace{0.1cm} SLA = Performance \vspace{0.15cm} Vendor Security Questions: "Are they compliant?" $\rightarrow$ Audit "Are they secure?" $\rightarrow$ Pen Test "Are they meeting service targets?" $\rightarrow$ SLA \end{block} \end{frame}
	% ================ ADDITIONAL DOMAIN 1 QUESTIONS ================
	
	\section{Domain 1: Additional Key Concepts}
	

	
	\begin{frame}{Q34: Risk Tolerance Example}
		\begin{block}{Question}
			Risk tolerance is best described as:
		\end{block}
		\begin{enumerate}[(A)]
			\item The maximum risk an enterprise can take
			\item The acceptable level of deviation from risk appetite
			\item The complete elimination of risk
			\item The transfer of risk to a third party
		\end{enumerate}
	\end{frame}
	
	\begin{frame}{Q34: Explanation}
		\begin{block}{Correct Answer: B}
			Risk tolerance is defined as the acceptable level of variation that management is willing to allow for any particular risk as the enterprise pursues its objectives.
		\end{block}
		\begin{block}{Option Analysis}
			\begin{itemize}
				\item \textbf{A - Incorrect:} This is risk capacity
				\item \textbf{B - Correct:} Tolerance is acceptable deviation
				\item \textbf{C - Incorrect:} Risk cannot be completely eliminated
				\item \textbf{D - Incorrect:} This is risk transfer
			\end{itemize}
		\end{block}
	\end{frame}
	
	
	\begin{frame}{Q34: Risk Tolerance} \begin{block}{Question} The board has approved a risk appetite that allows limited operational disruptions. Management establishes a threshold permitting customer-facing systems to be unavailable for no more than two hours per quarter before escalation is required. This threshold is BEST described as: \end{block} \begin{enumerate}[(A)] \item Risk capacity \item Risk tolerance \item Risk avoidance \item Risk transfer \end{enumerate} \begin{block}{CRISC Memory Aid} Risk Appetite = How much risk do we WANT to take? \vspace{0.1cm} Risk Tolerance = What limit triggers action? \vspace{0.1cm} Risk Capacity = How much risk can we SURVIVE? \vspace{0.15cm} Appetite Less than Capacity Tolerance = Measurable Threshold \end{block} 
	 \end{frame} \begin{frame}[shrink=10]{Q34: Explanation} \footnotesize \begin{block}{Correct Answer: B - Risk Tolerance} Risk tolerance defines the acceptable variation or threshold around a risk objective. It translates high-level risk appetite into measurable limits that management can monitor and enforce. \end{block} \begin{block}{Option Analysis} \textbf{A - Incorrect} Risk capacity represents the maximum amount of risk the organization can absorb before its viability or ability to achieve objectives is threatened. \vspace{0.15cm} \textbf{B - Correct} The two-hour threshold establishes the amount of deviation management is willing to accept before action is required. This is a classic example of risk tolerance. \vspace{0.15cm} \textbf{C - Incorrect} Risk avoidance involves eliminating the activity that creates the risk rather than establishing acceptable limits. \vspace{0.15cm} \textbf{D - Incorrect} Risk transfer shifts some or all of the financial impact of risk to another party, such as through insurance or contractual arrangements. \end{block}  \end{frame}
	
	\begin{frame}{Q35: Risk Culture Symptoms}
		\begin{block}{Question}
			Which is a symptom of an inadequate or problematic risk culture?
		\end{block}
		\begin{enumerate}[(A)]
			\item Open communication about risk
			\item Clear understanding of risk appetite
			\item Overconfidence in ability to manage risk
			\item Regular risk reporting to management
		\end{enumerate}
	\end{frame}
	
	\begin{frame}{Q35: Explanation}
		\begin{block}{Correct Answer: C}
			Overconfidence in ability to manage risk is a symptom of problematic risk culture.
		\end{block}
		\begin{block}{Option Analysis}
			\begin{itemize}
				\item \textbf{A - Incorrect:} Open communication is healthy
				\item \textbf{B - Incorrect:} Clear understanding is positive
				\item \textbf{C - Correct:} Overconfidence indicates problematic culture
				\item \textbf{D - Incorrect:} Regular reporting is good practice
			\end{itemize}
		\end{block}
	\end{frame}
	
	
	\begin{frame}{Q35: Risk Culture} \begin{block}{Question} During a review of the organization's risk management program, the board observes that business units frequently dismiss risk assessment findings because managers believe their teams can handle any issue that arises without additional controls. This is MOST indicative of: \end{block} \begin{enumerate}[(A)] \item Effective risk communication \item Strong understanding of risk appetite \item A problematic risk culture \item Mature risk reporting practices \end{enumerate} \end{frame} \begin{frame}[shrink=10]{Q35: Explanation} \footnotesize \begin{block}{Correct Answer: C - A Problematic Risk Culture} Excessive confidence in the organization's ability to manage risk can result in poor decision-making, inadequate controls, and a tendency to underestimate potential exposures. These are common indicators of an unhealthy risk culture. \end{block} \begin{block}{Option Analysis} \textbf{A - Incorrect} Effective risk communication encourages discussion, escalation, and transparency regarding risks. The scenario instead shows risk concerns being dismissed. \vspace{0.15cm} \textbf{B - Incorrect} Understanding risk appetite helps management make decisions within approved risk boundaries. Ignoring risk assessments does not demonstrate an understanding of risk appetite. \vspace{0.15cm} \textbf{C - Correct} Overconfidence, resistance to challenge, and dismissal of risk information are classic indicators of a weak risk culture. Organizations with these characteristics often underestimate risk exposure. \vspace{0.15cm} \textbf{D - Incorrect} Mature reporting practices ensure management receives accurate risk information. Receiving reports but ignoring them does not indicate reporting maturity. \end{block} \begin{block}{CRISC Memory Aid} Healthy Risk Culture: Open Communication Challenge Assumptions Escalate Concerns Understand Risk Appetite \vspace{0.15cm} Poor Risk Culture: Overconfidence Ignoring Warnings Blame Culture Fear of Escalation "It Won't Happen Here" \end{block} \end{frame}
	
	\begin{frame}{Q36: Confidentiality Breach}
		\begin{block}{Question}
			Which is an example of a breach of confidentiality?
		\end{block}
		\begin{enumerate}[(A)]
			\item Data being modified by unauthorized user
			\item Data being unavailable when needed
			\item Data being disclosed to unauthorized recipient
			\item Data being destroyed
		\end{enumerate}
	\end{frame}
	
	\begin{frame}{Q36: Explanation}
		\begin{block}{Correct Answer: C}
			A breach of confidentiality means the improper disclosure of information to a recipient not authorized access by the information owner.
		\end{block}
		\begin{block}{Option Analysis}
			\begin{itemize}
				\item \textbf{A - Incorrect:} This is a breach of integrity
				\item \textbf{B - Incorrect:} This is a breach of availability
				\item \textbf{C - Correct:} Unauthorized disclosure is confidentiality breach
				\item \textbf{D - Incorrect:} This is a breach of integrity/availability
			\end{itemize}
		\end{block}
	\end{frame}
	
	
	\begin{frame}{Q36: Confidentiality Breach}
		\begin{block}{Question}
			A payroll employee accidentally emails a file containing employee
			salaries and Social Security numbers to an external vendor that
			has no business need to access the information.
			
			Which information security objective has been MOST directly violated?
		\end{block}
		
		\begin{enumerate}[(A)]
			\item Integrity
			\item Availability
			\item Confidentiality
			\item Accountability
		\end{enumerate}
		
				\begin{block}{CRISC Memory Aid}
			
			Confidentiality
			
			= Unauthorized Disclosure
			
			\vspace{0.1cm}
			
			Integrity
			
			= Unauthorized Change
			
			\vspace{0.1cm}
			
			Availability
			
			= Unauthorized Interruption
			
			\vspace{0.1cm}
			
			Accountability
			
			= Who Did It?
			
		\end{block}
	\end{frame}
	
	\begin{frame}[shrink=10]{Q36: Explanation}
		\footnotesize
		
		\begin{block}{Correct Answer: C - Confidentiality}
			Confidentiality ensures that information is disclosed only to
			authorized individuals, processes, or systems. Sending sensitive
			information to an unauthorized recipient constitutes a breach of
			confidentiality.
		\end{block}
		
		\begin{block}{Option Analysis}
			
			\textbf{A - Incorrect: Integrity}
			
			Integrity relates to the accuracy, completeness, and
			trustworthiness of information. In the scenario, the data was
			disclosed but not altered.
			
			\vspace{0.15cm}
			
			\textbf{B - Incorrect: Availability}
			
			Availability ensures information is accessible when needed.
			The information remained available to authorized users.
			
			\vspace{0.15cm}
			
			\textbf{C - Correct: Confidentiality}
			
			The unauthorized disclosure of sensitive information is a
			classic confidentiality violation.
			
			\vspace{0.15cm}
			
			\textbf{D - Incorrect: Accountability}
			
			Accountability relates to tracing actions to specific users
			through logs, monitoring, and audit trails. Although
			accountability may help identify who sent the file, it is
			not the primary violation.
			
		\end{block}
		

		
	\end{frame}
	
	\begin{frame}{Q37: CIA Triad Trade-off}
		\begin{block}{Question}
			Increasing confidentiality typically:
		\end{block}
		\begin{enumerate}[(A)]
			\item Increases availability
			\item Decreases processing time
			\item Increases processing time and may reduce availability
			\item Has no effect on availability
		\end{enumerate}
	\end{frame}
	
	\begin{frame}{Q37: Explanation}
		\begin{block}{Correct Answer: C}
			Increasing confidentiality typically increases processing time, which reduces availability. There are trade-offs between the three principles.
		\end{block}
		\begin{block}{Option Analysis}
			\begin{itemize}
				\item \textbf{A - Incorrect:} Usually decreases availability
				\item \textbf{B - Incorrect:} Usually increases processing time
				\item \textbf{C - Correct:} Trade-off between confidentiality and availability
				\item \textbf{D - Incorrect:} There is typically an effect
			\end{itemize}
		\end{block}
	\end{frame}
	
	
	
	
	\begin{frame}[shrink=10]{Q37: Balancing Security and Business Needs}
		\begin{block}{Question}
			A financial institution plans to implement additional encryption
			and multifactor authentication controls to protect customer data.
			Business users are concerned that the new controls may slow access
			to critical applications.
			
			What is the MOST important consideration?
		\end{block}
		
		\begin{enumerate}[(A)]
			\item Confidentiality should always take precedence over availability
			\item Availability should always take precedence over confidentiality
			\item Security controls should be balanced with business requirements
			\item Additional controls should be avoided if they affect usability
		\end{enumerate}
		
		
				\begin{block}{CRISC Memory Aid}
			
			CIA = Balance
			
			Not Maximum Security
			
			Not Maximum Convenience
			
			\vspace{0.15cm}
			
			Business Objective
			
			$\rightarrow$ Risk Appetite
			
			$\rightarrow$ Security Controls
			
			$\rightarrow$ CIA Trade-offs
			
		\end{block}
		
	\end{frame}
	
	\begin{frame}[shrink=10]{Q37: Explanation}
		\footnotesize
		
		\begin{block}{Correct Answer: C - Security Controls Should Be Balanced with Business Requirements}
			Effective risk management requires balancing confidentiality,
			integrity, and availability to support business objectives.
			Increasing one objective may affect another, so management must
			determine the appropriate balance based on risk appetite and
			business needs.
		\end{block}
		
		\begin{block}{Option Analysis}
			
			\textbf{A - Incorrect}
			
			Confidentiality is important but does not automatically take
			priority over all other business requirements.
			
			\vspace{0.15cm}
			
			\textbf{B - Incorrect}
			
			Availability is critical for many services, but it should not
			override all confidentiality requirements.
			
			\vspace{0.15cm}
			
			\textbf{C - Correct}
			
			Risk management involves optimizing security while supporting
			business objectives. The goal is an appropriate balance rather
			than maximizing a single security objective.
			
			\vspace{0.15cm}
			
			\textbf{D - Incorrect}
			
			Usability is important, but controls should be evaluated based
			on risk rather than avoided because they introduce inconvenience.
			
		\end{block}
		

		
	\end{frame}
	
	% ================ ADDITIONAL DOMAIN 2 QUESTIONS ================
	
	\section{Domain 2: Additional Assessment Concepts}
	
	\begin{frame}{Q38: OCTAVE Methodology}
		\begin{block}{Question}
			OCTAVE focuses primarily on:
		\end{block}
		\begin{enumerate}[(A)]
			\item Technical vulnerabilities only
			\item Critical assets and risk to those assets
			\item Financial risk quantification
			\item Compliance requirements
		\end{enumerate}
	\end{frame}
	
	\begin{frame}{Q38: Explanation}
		\begin{block}{Correct Answer: B}
			OCTAVE focuses on critical assets and the risk to those assets using a comprehensive, systematic, context-driven and self-directed evaluation approach.
		\end{block}
		\begin{block}{Option Analysis}
			\begin{itemize}
				\item \textbf{A - Incorrect:} Not focused only on technical vulnerabilities
				\item \textbf{B - Correct:} Focuses on critical assets
				\item \textbf{C - Incorrect:} Not financial quantification
				\item \textbf{D - Incorrect:} Not primarily compliance
			\end{itemize}
		\end{block}
	\end{frame}
	
	
	
	
	\begin{frame}{Q38: OCTAVE Risk Assessment} \begin{block}{Question} A healthcare organization wants to perform a risk assessment that focuses on identifying its most critical information assets, understanding how those assets support business objectives, and evaluating risks in the context of the organization's operations. Which approach BEST supports this objective? \end{block} \begin{enumerate}[(A)] \item Focus primarily on technical vulnerabilities \item Focus on critical assets and risks to those assets \item Quantify financial loss for each identified threat \item Evaluate regulatory compliance requirements \end{enumerate} \begin{block}{CRISC Memory Aid} OCTAVE = Asset First = Business Context = Self-Directed Assessment \vspace{0.15cm} OCTAVE $\rightarrow$ Critical Assets FAIR $\rightarrow$ Financial Risk Vulnerability Scan $\rightarrow$ Technical Weaknesses Audit $\rightarrow$ Compliance \end{block} \end{frame} \begin{frame}[shrink=10]{Q38: Explanation} \footnotesize \begin{block}{Correct Answer: B - Focus on Critical Assets and Risks to Those Assets} OCTAVE (Operationally Critical Threat, Asset, and Vulnerability Evaluation) is an asset-centric risk assessment methodology. It focuses on identifying critical information assets and evaluating threats and vulnerabilities that could impact those assets. \end{block} \begin{block}{Option Analysis} \textbf{A - Incorrect} Vulnerability-focused approaches concentrate on technical weaknesses. OCTAVE starts with identifying critical assets and business impact rather than technical vulnerabilities. \vspace{0.15cm} \textbf{B - Correct} OCTAVE is built around understanding what assets are most important to the organization and assessing risks that could affect those assets. \vspace{0.15cm} \textbf{C - Incorrect} Financial quantification is more closely associated with quantitative risk assessment methodologies rather than OCTAVE's asset-centric approach. \vspace{0.15cm} \textbf{D - Incorrect} Compliance requirements may be considered during the assessment, but they are not the primary focus of OCTAVE. \end{block}  \end{frame}
	
	\begin{frame}{Q39: Phases of OCTAVE}
		\begin{block}{Question}
			Phase 1 of OCTAVE involves:
		\end{block}
		\begin{enumerate}[(A)]
			\item Identifying infrastructure vulnerabilities
			\item Building asset-based threat profiles
			\item Developing security strategy
			\item Implementing controls
		\end{enumerate}
	\end{frame}
	
	\begin{frame}{Q39: Explanation}
		\begin{block}{Correct Answer: B}
			Phase 1 of OCTAVE is to build asset-based threat profiles through organizational evaluation.
		\end{block}
		\begin{block}{Option Analysis}
			\begin{itemize}
				\item \textbf{A - Incorrect:} This is Phase 2
				\item \textbf{B - Correct:} Phase 1 builds threat profiles
				\item \textbf{C - Incorrect:} This is Phase 3
				\item \textbf{D - Incorrect:} This follows OCTAVE
			\end{itemize}
		\end{block}
	\end{frame}
	
	\begin{frame}{Q40: Fault Tree Analysis}
		\begin{block}{Question}
			Fault tree analysis is best described as:
		\end{block}
		\begin{enumerate}[(A)]
			\item Forward-looking, bottom-up model
			\item Top-down analysis starting with an event
			\item Statistical inference method
			\item Delphi consensus technique
		\end{enumerate}
	\end{frame}
	
	\begin{frame}{Q40: Explanation}
		\begin{block}{Correct Answer: B}
			A fault tree analysis starts with an event and examines possible means for the event to occur (top-down) and displays these results in a logical tree diagram.
		\end{block}
		\begin{block}{Option Analysis}
			\begin{itemize}
				\item \textbf{A - Incorrect:} This is event tree analysis
				\item \textbf{B - Correct:} Top-down analysis starting with event
				\item \textbf{C - Incorrect:} This is Bayesian analysis
				\item \textbf{D - Incorrect:} This is Delphi method
			\end{itemize}
		\end{block}
	\end{frame}
	
	
	
	\begin{frame}[shrink=10]{Q40: Root Cause Analysis}
		\begin{block}{Question}
			Following a major online banking outage, management wants to
			understand how multiple control failures, system errors, and
			human actions combined to cause the incident.
			
			Which risk analysis technique would BEST support this objective?
		\end{block}
		
		\begin{enumerate}[(A)]
			\item Event tree analysis
			\item Fault tree analysis
			\item Bayesian analysis
			\item Delphi technique
		\end{enumerate}
				
		\begin{block}{CRISC Memory Aid}
			
			Fault Tree
			
			= Start with FAILURE
			
			= Work Backward
			
			= Root Causes
			
			\vspace{0.15cm}
			
			Event Tree
			
			= Start with EVENT
			
			= Work Forward
			
			= Possible Outcomes
			
			\vspace{0.15cm}
			
			Delphi
			
			= Expert Opinions
			
			= Consensus
			
			\vspace{0.15cm}
			
			Bayesian
			
			= Probability Updates
			
		\end{block}
		
	\end{frame}
	
	\begin{frame}[shrink=10]{Q40: Explanation}
		\footnotesize
		
		\begin{block}{Correct Answer: B - Fault Tree Analysis}
			Fault Tree Analysis (FTA) is a top-down analytical technique
			used to determine the combinations of failures, control
			weaknesses, and events that could lead to a specific
			undesirable outcome.
		\end{block}
		
		\begin{block}{Option Analysis}
			
			\textbf{A - Incorrect}
			
			Event Tree Analysis (ETA) is a forward-looking technique that
			starts with an initiating event and examines possible outcomes.
			It asks: "What could happen next?"
			
			\vspace{0.15cm}
			
			\textbf{B - Correct}
			
			Fault Tree Analysis starts with an undesirable event and works
			backward to identify contributing causes and failures. It asks:
			"How did this happen?"
			
			\vspace{0.15cm}
			
			\textbf{C - Incorrect}
			
			Bayesian analysis uses probabilities and statistical inference
			to update likelihood estimates based on new information. It is
			not primarily a root-cause analysis technique.
			
			\vspace{0.15cm}
			
			\textbf{D - Incorrect}
			
			The Delphi technique obtains expert consensus through multiple
			rounds of anonymous feedback. It is used for estimation and
			forecasting rather than failure analysis.
			
		\end{block}

	\end{frame}
	
	\begin{frame}{Q41: Monte Carlo Analysis}
		\begin{block}{Question}
			Monte Carlo analysis is used to:
		\end{block}
		\begin{enumerate}[(A)]
			\item Eliminate risk entirely
			\item Establish aggregate variation in a system
			\item Identify single points of failure
			\item Document security policies
		\end{enumerate}
	\end{frame}
	
	\begin{frame}{Q41: Explanation}
		\begin{block}{Correct Answer: B}
			Monte Carlo simulation is used to establish the aggregate variation in a system resulting from variations in the system, for a number of inputs, where each input has a defined distribution.
		\end{block}
		\begin{block}{Option Analysis}
			\begin{itemize}
				\item \textbf{A - Incorrect:} Cannot eliminate risk entirely
				\item \textbf{B - Correct:} Establishes aggregate variation
				\item \textbf{C - Incorrect:} Not for identifying SPOFs
				\item \textbf{D - Incorrect:} Not for documentation
			\end{itemize}
		\end{block}
	\end{frame}
	
	\begin{frame}{Q42: Delphi Method}
		\begin{block}{Question}
			Delphi method leverages:
		\end{block}
		\begin{enumerate}[(A)]
			\item Expert opinion through multiple questionnaire rounds
			\item Statistical analysis of historical data
			\item Automated vulnerability scanning
			\item Document review
		\end{enumerate}
	\end{frame}
	
	\begin{frame}{Q42: Explanation}
		\begin{block}{Correct Answer: A}
			The Delphi method leverages expert opinion received using two or more rounds of questionnaires. After each round, results are summarized and communicated by a facilitator.
		\end{block}
		\begin{block}{Option Analysis}
			\begin{itemize}
				\item \textbf{A - Correct:} Expert opinion through multiple rounds
				\item \textbf{B - Incorrect:} Not statistical analysis
				\item \textbf{C - Incorrect:} Not automated scanning
				\item \textbf{D - Incorrect:} Not document review
			\end{itemize}
		\end{block}
	\end{frame}
	
	% ================ ADDITIONAL DOMAIN 3 QUESTIONS ================
	
	\section{Domain 3: Additional Response Concepts}
	
	\begin{frame}{Q43: Cost-Benefit Analysis Purpose}
		\begin{block}{Question}
			Cost-benefit analysis is used to:
		\end{block}
		\begin{enumerate}[(A)]
			\item Justify expense of control implementation
			\item Calculate annual loss expectancy
			\item Determine risk appetite
			\item Identify threats
		\end{enumerate}
	\end{frame}
	
	\begin{frame}{Q43: Explanation}
		\begin{block}{Correct Answer: A}
			Cost-benefit analysis is used to justify the expense associated with the implementation of controls. The expenditure on a control cannot be justified if the benefit realized is less than the cost.
		\end{block}
		\begin{block}{Option Analysis}
			\begin{itemize}
				\item \textbf{A - Correct:} Justifies control expense
				\item \textbf{B - Incorrect:} ALE is calculated differently
				\item \textbf{C - Incorrect:} Risk appetite is management decision
				\item \textbf{D - Incorrect:} Threats are identified through risk identification
			\end{itemize}
		\end{block}
	\end{frame}
	
	
	
	
	\begin{frame}{Q43: Cost-Benefit Analysis} \begin{block}{Question} A risk assessment identified a potential annual loss of \$500,000 from fraudulent transactions. Management is considering implementing a new fraud-detection solution that costs \$150,000 annually. Before approving the investment, management wants to determine whether the control is economically justified. Which analysis would BEST support this decision? \end{block} \begin{enumerate}[(A)] \item Cost-benefit analysis \item Annualized loss expectancy analysis \item Risk appetite assessment \item Threat assessment \end{enumerate}  \begin{block}{CRISC Memory Aid} Threat Assessment = What can happen? \vspace{0.1cm} ALE = How much can we lose? \vspace{0.1cm} Cost-Benefit Analysis = Is the control worth it? \vspace{0.1cm} Risk Appetite = How much risk will we accept? \vspace{0.15cm} Exam Formula: Risk First $\rightarrow$ Loss Estimate $\rightarrow$ Cost-Benefit $\rightarrow$ Control Decision \end{block} \end{frame}
	\begin{frame}[shrink=10]{Q43: Explanation} \footnotesize \begin{block}{Correct Answer: A - Cost-Benefit Analysis} Cost-benefit analysis compares the expected benefits of a control with its implementation and operating costs. It helps management determine whether the reduction in risk justifies the investment. \end{block} \begin{block}{Option Analysis} \textbf{A - Correct} Cost-benefit analysis is used when management must decide whether a proposed control provides sufficient value relative to its cost. \vspace{0.15cm} \textbf{B - Incorrect} Annualized Loss Expectancy (ALE) helps estimate expected losses. It provides important input but does not by itself determine whether a specific control investment is justified. \vspace{0.15cm} \textbf{C - Incorrect} Risk appetite defines how much risk the organization is willing to accept. It does not provide an economic comparison of control costs and benefits. \vspace{0.15cm} \textbf{D - Incorrect} Threat assessments identify potential threat sources and events. They do not determine whether a control is financially justified. \end{block}  \end{frame}
	
	\begin{frame}{Q44: Risk Transfer Example}
		\begin{block}{Question}
			Risk transfer is best exemplified by:
		\end{block}
		\begin{enumerate}[(A)]
			\item Purchasing insurance
			\item Implementing a firewall
			\item Conducting user training
			\item Creating a policy
		\end{enumerate}
	\end{frame}
	
	\begin{frame}{Q44: Explanation}
		\begin{block}{Correct Answer: A}
			Risk transfer is a decision to reduce loss by having another organization incur the cost. The most common example of risk transfer is the purchasing of insurance.
		\end{block}
		\begin{block}{Option Analysis}
			\begin{itemize}
				\item \textbf{A - Correct:} Insurance transfers risk
				\item \textbf{B - Incorrect:} Firewall is mitigation
				\item \textbf{C - Incorrect:} Training is mitigation/awareness
				\item \textbf{D - Incorrect:} Policy is administrative control
			\end{itemize}
		\end{block}
	\end{frame}
	
	

	
	\begin{frame}{Q48: Digital Signature}
		\begin{block}{Question}
			A digital signature provides:
		\end{block}
		\begin{enumerate}[(A)]
			\item Confidentiality only
			\item Integrity and nonrepudiation
			\item Availability only
			\item Authentication only
		\end{enumerate}
	\end{frame}
	
	\begin{frame}{Q48: Explanation}
		\begin{block}{Correct Answer: B}
			A digital signature combines a hash function with asymmetric encryption to verify the author's identity. It provides integrity (message not changed) and nonrepudiation (sender cannot deny sending).
		\end{block}
		\begin{block}{Option Analysis}
			\begin{itemize}
				\item \textbf{A - Incorrect:} Does not provide confidentiality on its own
				\item \textbf{B - Correct:} Provides integrity and nonrepudiation
				\item \textbf{C - Incorrect:} Not related to availability
				\item \textbf{D - Incorrect:} Provides more than authentication
			\end{itemize}
		\end{block}
	\end{frame}
	
	
	
	\begin{frame}[shrink=10]{Q48: Electronic Contract Approval}
		\begin{block}{Question}
			An organization requires vendors to electronically sign
			procurement contracts before they are processed. Management
			wants assurance that a signed contract has not been altered
			and that the signer cannot later deny signing it.
			
			Which security capability BEST meets this requirement?
		\end{block}
		
		\begin{enumerate}[(A)]
			\item Data encryption
			\item Digital signature
			\item Data backup
			\item Access control
		\end{enumerate}
				\begin{block}{CRISC Memory Aid}
			
			Encryption
			
			= Confidentiality
			
			\vspace{0.1cm}
			
			Hashing
			
			= Integrity
			
			\vspace{0.1cm}
			
			Digital Signature
			
			= Integrity + Authenticity + Nonrepudiation
			
			\vspace{0.1cm}
			
			Digital Signature
			
			$\neq$
			
			Confidentiality
			
			\vspace{0.15cm}
			
			Sign Then Encrypt
			
			= Confidentiality + Integrity + Nonrepudiation
			
		\end{block}
	\end{frame}
	
	\begin{frame}[shrink=10]{Q48: Explanation}
		\footnotesize
		
		\begin{block}{Correct Answer: B - Digital Signature}
			A digital signature uses hashing and asymmetric cryptography
			to verify the origin and integrity of information. It provides
			assurance that the message has not been altered and that the
			sender cannot reasonably deny having signed it.
		\end{block}
		
		\begin{block}{Option Analysis}
			
			\textbf{A - Incorrect}
			
			Encryption protects confidentiality by preventing
			unauthorized disclosure of information. Encryption alone
			does not prove who signed a document.
			
			\vspace{0.15cm}
			
			\textbf{B - Correct}
			
			Digital signatures provide:
			
			Integrity
			
			= Message was not altered.
			
			Authenticity
			
			= Signer identity can be verified.
			
			Nonrepudiation
			
			= Signer cannot deny the action.
			
			\vspace{0.15cm}
			
			\textbf{C - Incorrect}
			
			Backups support recovery and availability but do not
			validate document origin or integrity.
			
			\vspace{0.15cm}
			
			\textbf{D - Incorrect}
			
			Access controls restrict who may access a resource,
			but they do not prove that a specific user signed
			a document.
			
		\end{block}
		

		
	\end{frame}
	
	
	\begin{frame}[shrink=10]{Q49: Asymmetric Encryption Advantage}
		\begin{block}{Question}
			Asymmetric encryption's main advantage over symmetric is:
		\end{block}
		\begin{enumerate}[(A)]
			\item Faster processing speed
			\item Easier key management for secure exchange
			\item Lower computational requirements
			\item Simpler implementation
		\end{enumerate}
	\end{frame}
	
	\begin{frame}{Q49: Explanation}
		\begin{block}{Correct Answer: B}
			Asymmetric algorithms overcome the weakness of scalability of symmetric key cryptography by allowing secure key exchange. The public key can be freely distributed, making key management easier.
		\end{block}
		\begin{block}{Option Analysis}
			\begin{itemize}
				\item \textbf{A - Incorrect:} Asymmetric is slower
				\item \textbf{B - Correct:} Easier key management
				\item \textbf{C - Incorrect:} Higher computational requirements
				\item \textbf{D - Incorrect:} More complex implementation
			\end{itemize}
		\end{block}
	\end{frame}
	
	% ================ ADDITIONAL DOMAIN 4 QUESTIONS ================
	
	\section{Domain 4: Additional Monitoring Concepts}
	
	\begin{frame}{Q50: KRI Sensitivity}
		\begin{block}{Question}
			A KRI that generates too many false positives should be adjusted for:
		\end{block}
		\begin{enumerate}[(A)]
			\item Sensitivity
			\item Timeliness
			\item Frequency
			\item Reliability
		\end{enumerate}
	\end{frame}
	
	\begin{frame}{Q50: Explanation}
		\begin{block}{Correct Answer: A}
			KRIs that report on data points that cannot be controlled or are not alerting at the correct time should be adjusted (optimized) to be more precise, relevant or accurate. Sensitivity should be adjusted to reduce false positives.
		\end{block}
		\begin{block}{Option Analysis}
			\begin{itemize}
				\item \textbf{A - Correct:} Sensitivity affects false positives
				\item \textbf{B - Incorrect:} Timing is about when alerts occur
				\item \textbf{C - Incorrect:} Frequency is how often measured
				\item \textbf{D - Incorrect:} Reliability is correlation with risk
			\end{itemize}
		\end{block}
	\end{frame}
	
	
	
	
	\begin{frame}[shrink=10]{Q50: Optimizing KRIs}
		\begin{block}{Question}
			A risk practitioner implemented a Key Risk Indicator (KRI)
			to identify potential fraudulent transactions. After several
			months, management complains that the indicator generates a
			large number of alerts that require investigation but rarely
			identify actual risk events.
			
			What should the risk practitioner adjust FIRST?
		\end{block}
		
		\begin{enumerate}[(A)]
			\item Sensitivity thresholds
			\item Reporting timeliness
			\item Measurement frequency
			\item Data source reliability
		\end{enumerate}
				\begin{block}{CRISC Memory Aid}
			
			Too Many Alerts
			
			$\rightarrow$ Reduce Sensitivity
			
			\vspace{0.1cm}
			
			Missed Risk Events
			
			$\rightarrow$ Increase Sensitivity
			
			\vspace{0.1cm}
			
			Late Alerts
			
			$\rightarrow$ Improve Timeliness
			
			\vspace{0.1cm}
			
			Poor Data
			
			$\rightarrow$ Improve Reliability
			
		\end{block}
		
	\end{frame}
	
	\begin{frame}[shrink=10]{Q50: Explanation}
		\footnotesize
		
		\begin{block}{Correct Answer: A - Sensitivity Thresholds}
			KRIs should provide early warning of increasing risk exposure
			without generating excessive false positives. If an indicator
			triggers too frequently for events that do not represent
			meaningful risk, its sensitivity should be adjusted.
		\end{block}
		
		\begin{block}{Option Analysis}
			
			\textbf{A - Correct}
			
			Excessive false positives usually indicate that alert thresholds
			are too sensitive. Adjusting the sensitivity improves precision
			and reduces unnecessary investigations.
			
			\vspace{0.15cm}
			
			\textbf{B - Incorrect}
			
			Timeliness determines how quickly information is reported.
			Faster reporting does not reduce false positives.
			
			\vspace{0.15cm}
			
			\textbf{C - Incorrect}
			
			Frequency determines how often the KRI is measured or reported.
			Measuring less frequently may reduce the number of alerts but
			does not improve indicator accuracy.
			
			\vspace{0.15cm}
			
			\textbf{D - Incorrect}
			
			Data reliability is important, but the scenario indicates
			excessive alerts rather than inaccurate or inconsistent data.
			
		\end{block}
		

		
	\end{frame}
	
	\begin{frame}{Q51: KPI and KRI Relationship}
		\begin{block}{Question}
			KPIs and KRIs are often used together because:
		\end{block}
		\begin{enumerate}[(A)]
			\item They measure the same thing
			\item KPIs identify underperformance; KRIs provide early warnings
			\item KRIs replace KPIs in mature organizations
			\item They are independent and not related
		\end{enumerate}
	\end{frame}
	
	\begin{frame}{Q51: Explanation}
		\begin{block}{Correct Answer: B}
			KPIs help to identify underperforming aspects of the organization, while KRIs provide early warnings of increased risk. Failing to meet performance goals is a risk, so the two can be used together.
		\end{block}
		\begin{block}{Option Analysis}
			\begin{itemize}
				\item \textbf{A - Incorrect:} They measure different things
				\item \textbf{B - Correct:} Complementary purposes
				\item \textbf{C - Incorrect:} KRIs don't replace KPIs
				\item \textbf{D - Incorrect:} They are related
			\end{itemize}
		\end{block}
	\end{frame}
	
	
	\begin{frame}{Q51: KPI and KRI Usage}
		\begin{block}{Question}
			Executive management reviews a dashboard showing:
			
			\begin{itemize}
				\item Mean incident response time
				\item Percentage of systems patched on schedule
				\item Increase in privileged-access violations
				\item Growth in phishing-related events
			\end{itemize}
			
			Why is it beneficial to monitor both KPIs and KRIs together?
		\end{block}
		
		\begin{enumerate}[(A)]
			\item They measure the same organizational objectives
			\item Performance metrics and risk indicators provide complementary information
			\item KRIs eliminate the need for performance measurement
			\item KPIs are useful only when KRIs exceed thresholds
		\end{enumerate}
	\end{frame}
	
	\begin{frame}[shrink=10]{Q51: Explanation}
		\footnotesize
		
		\begin{block}{Correct Answer: B - Performance Metrics and Risk Indicators Provide Complementary Information}
			KPIs measure whether objectives are being achieved, while
			KRIs provide insight into increasing risk exposure. Together,
			they help management understand both organizational performance
			and emerging threats that could affect future performance.
		\end{block}
		
		\begin{block}{Option Analysis}
			
			\textbf{A - Incorrect}
			
			KPIs and KRIs are related but measure different concepts.
			KPIs focus on outcomes and performance, whereas KRIs focus
			on risk exposure.
			
			\vspace{0.15cm}
			
			\textbf{B - Correct}
			
			An organization can appear to be performing well while risk
			exposure is increasing. Monitoring both performance and risk
			provides a more complete management view.
			
			\vspace{0.15cm}
			
			\textbf{C - Incorrect}
			
			KRIs supplement KPIs; they do not replace them.
			
			\vspace{0.15cm}
			
			\textbf{D - Incorrect}
			
			KPIs remain valuable regardless of KRI thresholds because
			performance measurement is a continuous management activity.
			
		\end{block}
		
		\begin{block}{CRISC Memory Aid}
			
			KPI= Are We Achieving Goals?
			
			\vspace{0.1cm}
			
			KRI			= Are Risks Increasing?
			
			\vspace{0.1cm}
			
			KCI			= Are Controls Working?
			
			\vspace{0.15cm}
			
			Response Time			= KPI
			
			\vspace{0.1cm}
			
			Increase in Security Violations			= KRI
			
			\vspace{0.1cm}
			
			Patch Compliance			= KCI
			
		\end{block}
		
	\end{frame}
	
	
	
	
	\begin{frame}{Q52: KRI Maintenance Frequency}
		\begin{block}{Question}
			KRIs should be evaluated:
		\end{block}
		\begin{enumerate}[(A)]
			\item Annually
			\item Only when incidents occur
			\item On a regular basis
			\item Never after implementation
		\end{enumerate}
	\end{frame}
	
	\begin{frame}{Q52: Explanation}
		\begin{block}{Correct Answer: C}
			Because the organization's internal and external environments are constantly changing, the set of KRIs needs to be evaluated on a regular basis to verify that each KRI remains properly related to the risk appetite and tolerance levels.
		\end{block}
		\begin{block}{Option Analysis}
			\begin{itemize}
				\item \textbf{A - Incorrect:} May need more frequent evaluation
				\item \textbf{B - Incorrect:} Should be proactive, not reactive
				\item \textbf{C - Correct:} Regular evaluation is needed
				\item \textbf{D - Incorrect:} Must be maintained over time
			\end{itemize}
		\end{block}
	\end{frame}
	
	\begin{frame}{Q53: SIEM Primary Function}
		\begin{block}{Question}
			SIEM systems primarily:
		\end{block}
		\begin{enumerate}[(A)]
			\item Correlate data from multiple sources
			\item Store backup data
			\item Manage user identities
			\item Encrypt network traffic
		\end{enumerate}
	\end{frame}
	
	\begin{frame}{Q53: Explanation}
		\begin{block}{Correct Answer: A}
			SIEM systems are integrated data correlation tools that capture data from multiple sources and analyze system, application and network activity for possible security events.
		\end{block}
		\begin{block}{Option Analysis}
			\begin{itemize}
				\item \textbf{A - Correct:} Correlation of multiple data sources
				\item \textbf{B - Incorrect:} Not for backup storage
				\item \textbf{C - Incorrect:} Identity management is separate
				\item \textbf{D - Incorrect:} Not for traffic encryption
			\end{itemize}
		\end{block}
	\end{frame}
	
	\begin{frame}{Q54: Log Review Challenges}
		\begin{block}{Question}
			Common challenge with log review is:
		\end{block}
		\begin{enumerate}[(A)]
			\item Too much data to analyze effectively
			\item Not enough data being collected
			\item Logs being too expensive
			\item Logs not being required
		\end{enumerate}
	\end{frame}
	
	\begin{frame}{Q54: Explanation}
		\begin{block}{Correct Answer: A}
			Common challenges relating to the effective use of logs include having too much data, difficulty in searching for relevant information, improper configuration, and modification or deletion of data before it is read.
		\end{block}
		\begin{block}{Option Analysis}
			\begin{itemize}
				\item \textbf{A - Correct:} Data volume is a challenge
				\item \textbf{B - Incorrect:} Often there is too much data
				\item \textbf{C - Incorrect:} Cost is not the primary challenge
				\item \textbf{D - Incorrect:} Logs are typically required
			\end{itemize}
		\end{block}
	\end{frame}
	
	
	
	\begin{frame}[shrink=10]{Q54: Security Monitoring}
		\begin{block}{Question}
			A financial institution collects security logs from firewalls,
			databases, application servers, and endpoint systems. Although
			the organization gathers millions of events daily, several
			significant incidents were not identified until weeks later.
			
			What is the MOST likely reason?
		\end{block}
		
		\begin{enumerate}[(A)]
			\item Excessive log volume makes effective analysis difficult
			\item Insufficient logging is occurring
			\item Log storage costs are too high
			\item Log collection is not required
		\end{enumerate}
				\begin{block}{CRISC Memory Aid}
			
			Collecting Logs
			
			$\neq$
			
			Monitoring Logs
			
			$\neq$
			
			Analyzing Logs
			
			\vspace{0.15cm}
			
			Log Challenge
			
			= Too Much Data
			
			\vspace{0.1cm}
			
			SIEM
			
			= Correlate Data
			
			= Reduce Noise
			
			= Highlight Risk
			
			\vspace{0.1cm}
			
			Common Exam Rule:
			
			Data Exists
			
			but
			
			Actionable Information Is Missing
			
		\end{block}
	\end{frame}
	
	\begin{frame}[shrink=10]{Q54: Explanation}
		\footnotesize
		
		\begin{block}{Correct Answer: A - Excessive Log Volume Makes Effective Analysis Difficult}
			Modern organizations often generate enormous quantities of
			log data. The primary challenge is rarely data collection;
			it is identifying meaningful events and anomalies within the
			volume of information collected.
		\end{block}
		
		\begin{block}{Option Analysis}
			
			\textbf{A - Correct}
			
			Security teams frequently struggle with information overload.
			Large volumes of logs can make it difficult to identify
			critical events, suspicious patterns, and emerging threats.
			
			\vspace{0.15cm}
			
			\textbf{B - Incorrect}
			
			The scenario indicates that millions of events are already
			being collected. The issue is analysis, not collection.
			
			\vspace{0.15cm}
			
			\textbf{C - Incorrect}
			
			Storage costs may be a concern, but they would not explain
			why incidents remained undetected despite available log data.
			
			\vspace{0.15cm}
			
			\textbf{D - Incorrect}
			
			Security logs are commonly required for monitoring,
			investigations, compliance, and incident response.
			
		\end{block}
		

		
	\end{frame}
	
	\begin{frame}{Q55: Control Monitoring Purpose}
		\begin{block}{Question}
			Purpose of control monitoring is to verify:
		\end{block}
		\begin{enumerate}[(A)]
			\item Whether the control is working
			\item Whether the control is effectively addressing the risk
			\item Whether staff is following procedures
			\item Whether logs are being reviewed
		\end{enumerate}
	\end{frame}
	
	\begin{frame}{Q55: Explanation}
		\begin{block}{Correct Answer: B}
			Because the purpose of a control is to mitigate a risk, the purpose of control monitoring is to verify whether the control is effectively addressing the risk, not to see whether the control itself is working.
		\end{block}
		\begin{block}{Option Analysis}
			\begin{itemize}
				\item \textbf{A - Incorrect:} This is just operational status
				\item \textbf{B - Correct:} Verify effectiveness against risk
				\item \textbf{C - Incorrect:} This is compliance monitoring
				\item \textbf{D - Incorrect:} This is log review, not control monitoring
			\end{itemize}
		\end{block}
	\end{frame}
	
	
	
	\begin{frame}[shrink=10]{Q55: Control Monitoring}
		\begin{block}{Question}
			A financial institution implemented multifactor authentication
			to reduce the risk of unauthorized access to customer accounts.
			Six months later, management asks for evidence that the control
			is providing value.
			
			What should control monitoring primarily determine?
		\end{block}
		
		\begin{enumerate}[(A)]
			\item Whether the multifactor authentication system is operating correctly
			\item Whether the control is reducing unauthorized-access risk to an acceptable level
			\item Whether employees are complying with security procedures
			\item Whether authentication logs are being reviewed
		\end{enumerate}
				
		\begin{block}{CRISC Memory Aid}
			
			Control Working?	
			= Operating Effectiveness
			
			\vspace{0.1cm}
			
			Risk Reduced?
						= Control Effectiveness
			
			\vspace{0.15cm}
			
			CRISC usually cares more about:
						``Did risk decrease?''
						than			
			``Did the control run?''
			
		\end{block}
	\end{frame}
	
	\begin{frame}[shrink=10]{Q55: Explanation}
		\footnotesize
		
		\begin{block}{Correct Answer: B - Whether the Control Is Reducing Risk}
			Controls exist to manage risk, not merely to operate correctly.
			The ultimate measure of control effectiveness is whether the
			control reduces risk exposure to a level consistent with the
			organization's risk appetite and tolerance.
		\end{block}
		
		\begin{block}{Option Analysis}
			
			\textbf{A - Incorrect}
			
			A control may function exactly as designed yet still fail to
			adequately mitigate the target risk. Operational status alone
			does not prove effectiveness.
			
			\vspace{0.15cm}
			
			\textbf{B - Correct}
			
			The purpose of control monitoring is to determine whether the
			control is effectively reducing the associated risk and helping
			achieve the intended control objective.
			
			\vspace{0.15cm}
			
			\textbf{C - Incorrect}
			
			Procedure compliance may support control effectiveness, but
			compliance monitoring is not the primary objective of control
			monitoring.
			
			\vspace{0.15cm}
			
			\textbf{D - Incorrect}
			
			Log reviews help assess operational activity but do not directly
			demonstrate whether the underlying risk is being reduced.
			
		\end{block}

		
	\end{frame}
	% ================ FINAL SLIDES ================
	
	\begin{frame}{CRISC Key Takeaways from Review Manual}
		\begin{block}{Essential Concepts}
			\begin{itemize}
				\item Risk appetite = amount of risk willing to accept in pursuit of mission
				\item Risk tolerance = acceptable level of deviation from appetite
				\item Risk capacity = maximum risk before existence is questioned
				\item C MM: 5 levels from Incomplete to Optimized
				\item CIA: Confidentiality, Integrity, Availability
				\item Digital signatures provide integrity and nonrepudiation
			\end{itemize}
		\end{block}
	\end{frame}
	
	\begin{frame}{CRISC Key Takeaways (continued)}
		\begin{block}{More Essential Concepts}
			\begin{itemize}
				\item Risk assessment = identify and evaluate risk and potential effects
				\item Risk response options: Accept, Mitigate, Transfer, Avoid
				\item Residual risk = Inherent risk - Controls
				\item KRIs predict risk; KPIs measure performance
				\item 3 LoD: Operational management, Risk oversight, Audit
			\end{itemize}
		\end{block}
	\end{frame}
	
	\begin{frame}{Final Advice}
		\begin{block}{Exam Preparation}
			\centering
			\vspace{0.3cm}
			\textbf{Remember:}
			\begin{itemize}
				\item Understand the \textit{why} behind each concept
				\item Know the difference between risk appetite, tolerance, and capacity
				\item Understand the risk management life cycle (Identify, Assess, Respond, Monitor)
				\item Practice scenario-based questions
				\item Use the ISACA mindset - choose the MOST appropriate answer
			\end{itemize}
			\vspace{0.5cm}
			\textbf{Good Luck on Your CRISC Exam!}
		\end{block}
	\end{frame}
	
	
	\begin{frame}[shrink=5]{CRISC Ultimate Memory Sheet - Part 1}
		\footnotesize
		
		\textbf{Risk Hierarchy (Most Tested)}
		
		Risk Appetite
		= Amount of risk management WANTS to take.
		
		Risk Tolerance
		= Allowed variation from appetite.
		
		Risk Capacity
		= Maximum risk organization can SURVIVE.
		
		\medskip
		
		Think:
		
		Appetite = Desire
		
		Tolerance = Limit
		
		Capacity = Breaking Point
		
		\medskip
		
		Example:
		
		"We want aggressive growth"
		= Appetite
		
		"Losses up to \$5M acceptable"
		= Tolerance
		
		"Losses above \$50M threaten company"
		= Capacity
		
		\medskip
		
		Rule:
		
		Appetite Less Than Capacity
		
	\end{frame}
	
	\begin{frame}[shrink=5]{CRISC Ultimate Memory Sheet - Part 2}
		\footnotesize
		
		\textbf{Business Continuity}
		
		BIA
		= Business Impact Analysis
		
		BIA Outputs:
		
		RTO
		= Recovery Time Objective
		
		RPO
		= Recovery Point Objective
		
		MTD
		= Maximum Tolerable Downtime
		
		MAO
		= Maximum Acceptable Outage
		
		\medskip
		
		Memory:
		
		RPO = Data Loss
		
		RTO = Downtime
		
		MTD = Survival
		
		\medskip
		
		Exam Rule:
		
		RTO must ALWAYS be less than or equal to MTD
		
	\end{frame}
	
	\begin{frame}[shrink=5]{CRISC Ultimate Memory Sheet - Part 3}
		\footnotesize
		
		\textbf{Control Types}
		
		Preventive
		= Stop it
		
		Detective
		= Find it
		
		Corrective
		= Fix it
		
		Recovery
		= Restore it
		
		Deterrent
		= Discourage it
		
		Compensating
		= Alternative control
		
		Directive
		= Tell people what to do
		
		\medskip
		
		Examples:
		
		MFA = Preventive
		
		Firewall = Preventive
		
		IDS = Detective
		
		SIEM = Detective
		
		Backup Restore = Corrective
		
		Warning Banner = Deterrent
		
	\end{frame}
	
	\begin{frame}[shrink=5]{CRISC Ultimate Memory Sheet - Part 4}
		\footnotesize
		
		\textbf{Ownership Traps}
		
		Business Owns Risk
		
		Management Accepts Risk
		
		IT Implements Controls
		
		Risk Management Facilitates Process
		
		Audit Provides Assurance
		
		\medskip
		
		Most Common Exam Trap:
		
		Who owns risk?
		
		NOT IT
		
		NOT Security
		
		NOT Audit
		
		Answer:
		
		Business Management
		
		\medskip
		
		Responsibility can be delegated.
		
		Accountability cannot.
		
	\end{frame}
	
	\begin{frame}[shrink=5]{CRISC Ultimate Memory Sheet - Part 5}
		\footnotesize
		
		\textbf{CIA Triad}
		
		Confidentiality
		= Unauthorized Disclosure
		
		Integrity
		= Unauthorized Modification
		
		Availability
		= Unauthorized Interruption
		
		\medskip
		
		Examples:
		
		Data emailed to wrong person
		= Confidentiality
		
		Database changed without approval
		= Integrity
		
		Application outage
		= Availability
		
		\medskip
		
		Weather Forecast System
		
		Most Important = Integrity
		
		Hospital Emergency System
		
		Most Important = Availability
		
	\end{frame}
	
	\begin{frame}[shrink=5]{CRISC Ultimate Memory Sheet - Part 6}
		\footnotesize
		
		\textbf{Cryptography}
		
		Encryption
		= Confidentiality
		
		Hashing
		= Integrity
		
		Digital Signature
		= Integrity + Authenticity + Nonrepudiation
		
		PKI
		= Public Key Infrastructure
		
		CA
		= Certificate Authority
		
		CRL
		= Certificate Revocation List
		
		CSR
		= Certificate Signing Request
		
		\medskip
		
		Exam Trap:
		
		Digital Signature
		
		DOES NOT provide confidentiality
		
		Hashing
		
		IS NOT Encryption
		
	\end{frame}
	
	\begin{frame}[shrink=5]{CRISC Ultimate Memory Sheet - Part 7}
		\footnotesize
		
		\textbf{Risk Response}
		
		Accept
		= Keep risk
		
		Mitigate
		= Add controls
		
		Transfer
		= Insurance / Outsourcing
		
		Avoid
		= Stop activity
		
		\medskip
		
		Questions usually look like:
		
		Buy insurance
		= Transfer
		
		Implement MFA
		= Mitigate
		
		Accept residual risk
		= Accept
		
		Cancel project
		= Avoid
		
	\end{frame}
	
	\begin{frame}[shrink=5]{CRISC Ultimate Memory Sheet - Part 8}
		\footnotesize
		
		\textbf{KPI vs KRI vs KCI}
		
		KPI
		= Are we achieving objectives?
		
		KRI
		= Is risk increasing?
		
		KCI
		= Are controls working?
		
		\medskip
		
		Examples:
		
		Customer Satisfaction
		= KPI
		
		Incident Response Time
		= KPI
		
		Failed Logins
		= KRI
		
		Privileged Access Violations
		= KRI
		
		Patch Compliance
		= KCI
		
		MFA Coverage
		= KCI
		
	\end{frame}
	
	\begin{frame}[shrink=5]{CRISC Ultimate Memory Sheet - Part 9}
		\footnotesize
		
		\textbf{Risk Calculations}
		
		Asset
		$\rightarrow$
		Threat
		$\rightarrow$
		Vulnerability
		$\rightarrow$
		Inherent Risk
		$\rightarrow$
		Controls
		$\rightarrow$
		Residual Risk
		
		\medskip
		
		Inherent Risk
		
		= Before Controls
		
		Residual Risk
		
		= After Controls
		
		\medskip
		
		Strong Controls
		
		$\rightarrow$
		Lower Residual Risk
		
		Weak Controls
		
		$\rightarrow$
		Higher Residual Risk
		
	\end{frame}
	
	\begin{frame}[shrink=5]{CRISC Final 60-Second Brain Dump}
		\footnotesize
		
		Business Owns Risk
		
		Management Accepts Risk
		
		Audit Provides Assurance
		
		\medskip
		
		BIA $\rightarrow$ RTO, RPO, MTD
		
		CMM $\rightarrow$ Maturity
		
		Gap Analysis $\rightarrow$ Control Deficiencies
		
		Cost-Benefit $\rightarrow$ Justify Controls
		
		\medskip
		
		KPI = Performance
		
		KRI = Risk
		
		KCI = Controls
		
		\medskip
		
		Encryption = Confidentiality
		
		Hashing = Integrity
		
		Signature = Nonrepudiation
		
		\medskip
		
		Accept = Keep
		
		Mitigate = Control
		
		Transfer = Insurance
		
		Avoid = Stop Activity
		
		\medskip
		
		CRISC Golden Rule:
		
		FIRST identify.
		
		THEN assess.
		
		THEN analyze.
		
		THEN respond.
		
		THEN monitor.
		
		Never jump directly to implementation.
		
	\end{frame}
	
\end{document}